%% Volume II, Chapter 8: Symmetric Matrices and Quadratic Forms
\chapter{Symmetric Matrices and Quadratic Forms}
\label{ch:symmetric-matrices}

\abstract{Symmetric matrices are the matrices met most often in scientific computing. The Hessian of a function, the covariance matrix of data, the Gram matrix of a least squares problem and the stiffness matrix of a structure are all symmetric, as is every matrix that describes an energy or a variance. For such matrices the eigenvalues are real, the eigenvectors can be chosen at right angles to one another, and in their coordinates every quadratic form becomes a plain sum of squares. Whether a critical point is a minimum, what shape an ellipsoid has and how well a system is conditioned can then be read off from the eigenvalues, which, unlike those of general matrices, hardly move when the entries are perturbed.}

Symmetric matrices have appeared in almost every chapter of this part, and many questions about them were left open. \Vcref{1}{sec:extrema-n} classified critical points by the sign behaviour of the Hessian and stated two criteria for definiteness, by eigenvalues and by leading principal minors, in \vcref{1}{thm:definiteness-criteria}, with the proofs deferred; \vcref{1}{sec:lagrange} found that the extreme values of a quadratic form on the unit sphere are eigenvalues. \Vcref{1}{thm:gaussian-quadratic} evaluated the Gaussian integral of a positive definite quadratic form completely only for two variables. \Cref{ch:vectors-matrices} showed that Gram matrices are positive semidefinite and announced that the Schur complement of a positive definite matrix is positive definite. \Cref{ch:linear-systems} proved that a symmetric matrix with a factorisation $LDL^{\mathsf T}$ with positive pivots is positive definite and announced the converse and the Cholesky factorisation, and \cref{ch:determinants} expressed the pivots by leading principal minors. \Cref{ch:inner-products} found that every inner product on $\mathbb F^n$ has the form $\mathbf x^{\mathsf H}G\mathbf y$ and that the QR factorisation produces a factorisation $A^{\mathsf H}A=R^{\mathsf H}R$, and \cref{ch:eigenvalues} proved that the eigenvalues of a Hermitian matrix are real and that its eigenvectors for different eigenvalues are orthogonal, but left open whether a Hermitian matrix always has enough eigenvectors. This chapter answers all these questions.

The central result is the spectral theorem: every Hermitian matrix $A$ can be written as $A=Q\Lambda Q^{\mathsf H}$ with a unitary matrix $Q$ and a real diagonal matrix $\Lambda$, so that $\mathbb F^n$ has an orthonormal basis of eigenvectors. In the coordinates of this basis the quadratic form $\mathbf x^{\mathsf H}A\mathbf x$ is a weighted sum of squares, and its sign behaviour, its level sets and its extreme values on the sphere are read off from the eigenvalues. The positive definite matrices are characterised in seven equivalent ways, by eigenvalues, by factorisations, by minors and by pivots, and the last of these leads to the Cholesky factorisation $A=R^{\mathsf H}R$, which solves symmetric positive definite systems at half the cost of Gaussian elimination and is at the same time the cheapest test for positive definiteness. The Rayleigh quotient characterises the largest and the smallest eigenvalue by an optimisation problem, and the Courant--Fischer theorem extends this to every eigenvalue; its consequences, the inequalities of Weyl and the interlacing theorem of Cauchy, show that the eigenvalues of Hermitian matrices are perfectly conditioned. Sylvester's law of inertia describes what a change of variables can and cannot do to a quadratic form, and it turns the pivots of elimination into a way of counting eigenvalues. The chapter ends with the Gaussian integral in $n$ variables and other uses, and with experiments.

As in the preceding chapters, $\mathbb F$ denotes $\mathbb R$ or $\mathbb C$. We state the results for Hermitian matrices, $A^{\mathsf H}=A$; a real symmetric matrix is a Hermitian matrix with real entries, and for it every statement holds with real vectors in place of complex ones and orthogonal matrices in place of unitary ones. Where the real case needs an argument of its own, we give it. The eigenvalues of a Hermitian matrix $A\in\mathbb F^{n\times n}$ are real, and throughout the chapter they are numbered in \emph{increasing} order and repeated according to their algebraic multiplicity,
\[
\lambda_1(A)\le\lambda_2(A)\le\cdots\le\lambda_n(A),\qquad\lambda_{\min}(A)\coloneqq\lambda_1(A),\quad\lambda_{\max}(A)\coloneqq\lambda_n(A);
\]
we write $\lambda_k$ for $\lambda_k(A)$ when the matrix is clear. This is the order in which \cref{thm:second-difference-eigen} listed the eigenvalues of $T_n$. The power method of \cref{sec:computing-eigenvalues} numbered the eigenvalues of a general matrix by decreasing modulus, a different convention for a different purpose, which we use only when we return to it, and the singular values of \cref{ch:svd} are numbered in decreasing order, as is customary. The singular value decomposition of \cref{ch:svd} and the Schur and Jordan forms of \cref{ch:canonical-forms} are not used in this chapter.

The chapter is organised as follows. \Cref{sec:quadratic-forms} introduces quadratic and Hermitian forms, their matrices and the change of variables. \Cref{sec:spectral-theorem} proves the spectral theorem, for Hermitian and for normal matrices and for commuting families, and \cref{sec:principal-axes} applies it to quadratic forms, conics, ellipsoids and quadratic functions. \Cref{sec:positive-definite} characterises positive definite and semidefinite matrices and collects their properties, and \cref{sec:cholesky} treats the Cholesky factorisation, the square root of a positive semidefinite matrix and weighted least squares. \Cref{sec:rayleigh} studies the Rayleigh quotient, \cref{sec:min-max} the min--max principles and their consequences, and \cref{sec:inertia} the law of inertia, indefinite forms and the generalised eigenvalue problem. \Cref{sec:gaussian-uses} evaluates Gaussian integrals and collects other uses, and \cref{sec:symmetric-computation} looks at computation.

\section{Quadratic Forms}
\label{sec:quadratic-forms}

\subsection{Forms and Their Matrices}

A \term{quadratic form} in $n$ real variables is a homogeneous polynomial of degree $2$,
\[
q(\mathbf x)=\sum_{i=1}^n\sum_{j=1}^na_{ij}x_ix_j=\langle\mathbf x,A\mathbf x\rangle ,
\]
with a matrix $A=(a_{ij})\in\mathbb R^{n\times n}$. Each product $x_ix_j$ with $i\ne j$ occurs twice in the double sum, with the coefficients $a_{ij}$ and $a_{ji}$, so only their sum matters, and by \cref{prop:symmetric-parts} the form does not change when $A$ is replaced by its symmetric part $(A+A^{\mathsf T})/2$. We therefore take $A$ symmetric. Then the coefficient of $x_i^2$ in $q$ is $a_{ii}$ and the coefficient of $x_ix_j$, $i<j$, is $2a_{ij}$, so the matrix can be read off from the polynomial. For example,
\[
2x_1^2+2x_1x_2+2x_2^2=\Bigl\langle\mathbf x,\begin{bmatrix}2&1\\1&2\end{bmatrix}\mathbf x\Bigr\rangle,\qquad x_1^2-x_3^2+4x_1x_2-6x_2x_3=\Bigl\langle\mathbf x,\begin{bmatrix}1&2&0\\2&0&-3\\0&-3&-1\end{bmatrix}\mathbf x\Bigr\rangle .
\]
Over $\mathbb C$ the inner product $\langle\mathbf x,\mathbf y\rangle=\mathbf x^{\mathsf H}\mathbf y$ is conjugate linear in its first argument, and the natural counterpart of a quadratic form is $\mathbf x^{\mathsf H}A\mathbf x$.

\begin{definition}[Quadratic, sesquilinear and Hermitian forms]
\label{def:quadratic-form}
Let $A\in\mathbb F^{n\times n}$. The function $b_A(\mathbf x,\mathbf y)\coloneqq\langle\mathbf x,A\mathbf y\rangle=\mathbf x^{\mathsf H}A\mathbf y$ on $\mathbb F^n\times\mathbb F^n$ is the \term{sesquilinear form} of $A$, \term{bilinear} if $\mathbb F=\mathbb R$, and $q_A(\mathbf x)\coloneqq\langle\mathbf x,A\mathbf x\rangle=\mathbf x^{\mathsf H}A\mathbf x$ is its \term{quadratic form}. For $\mathbb F=\mathbb C$ and Hermitian $A$, $q_A$ is called a \term{Hermitian form}.
\end{definition}

The form $b_A$ is linear in $\mathbf y$ and conjugate linear in $\mathbf x$, by \cref{prop:hermitian-inner-product}\,(ii); ``sesquilinear'' means ``one and a half times linear''. A quadratic form satisfies $q_A(t\mathbf x)=|t|^2q_A(\mathbf x)$ for every scalar $t$, so its values on the unit sphere determine it everywhere. The next proposition shows that a form determines its matrix, and that over $\mathbb C$ the values of a form are real exactly for Hermitian matrices.

\begin{proposition}[Forms and matrices]
\label{prop:form-matrix}
\leavevmode
\begin{enumerate}
\item (Polarisation) For $A\in\mathbb C^{n\times n}$ and $\mathbf x,\mathbf y\in\mathbb C^n$,
\[
\mathbf x^{\mathsf H}A\mathbf y=\frac14\Bigl(q_A(\mathbf x+\mathbf y)-q_A(\mathbf x-\mathbf y)-i\,q_A(\mathbf x+i\mathbf y)+i\,q_A(\mathbf x-i\mathbf y)\Bigr),
\]
and for a real symmetric $A$ and $\mathbf x,\mathbf y\in\mathbb R^n$, $\langle\mathbf x,A\mathbf y\rangle=\bigl(q_A(\mathbf x+\mathbf y)-q_A(\mathbf x-\mathbf y)\bigr)/4$.
\item Two real symmetric matrices with the same quadratic form on $\mathbb R^n$ are equal. Two complex matrices with the same quadratic form on $\mathbb C^n$ are equal, symmetric or not.
\item For $A\in\mathbb C^{n\times n}$, the number $q_A(\mathbf x)$ is real for every $\mathbf x\in\mathbb C^n$ if and only if $A$ is Hermitian.
\end{enumerate}
\end{proposition}

\begin{proof}
(i) For a scalar $t$, the rules of \cref{prop:hermitian-inner-product}\,(ii) give
\[
q_A(\mathbf x+t\mathbf y)=q_A(\mathbf x)+|t|^2q_A(\mathbf y)+t\,\mathbf x^{\mathsf H}A\mathbf y+\bar t\,\mathbf y^{\mathsf H}A\mathbf x .
\]
Multiply by $\bar t$ and add the four equations for $t=1,-1,i,-i$. The sums of $\bar t$, of $\bar t|t|^2$ and of $\bar t^2$ over these four numbers are $0$, while the sum of $\bar tt=1$ is $4$, so only $4\,\mathbf x^{\mathsf H}A\mathbf y$ remains; the weights $\bar t$ are $1,-1,-i,i$, which gives the formula. In the real symmetric case, $\mathbf y^{\mathsf T}A\mathbf x=\mathbf x^{\mathsf T}A\mathbf y$, and subtracting the equations for $t=-1$ and $t=1$ gives $4\langle\mathbf x,A\mathbf y\rangle$. (ii) By (i), the form determines $\langle\mathbf x,A\mathbf y\rangle$ for all $\mathbf x$ and $\mathbf y$, in particular $a_{ij}=\langle\mathbf e_i,A\mathbf e_j\rangle$. (iii) For every $A$ and $\mathbf x$, $\overline{q_A(\mathbf x)}=(\mathbf x^{\mathsf H}A\mathbf x)^{\mathsf H}=\mathbf x^{\mathsf H}A^{\mathsf H}\mathbf x=q_{A^{\mathsf H}}(\mathbf x)$. So if $A^{\mathsf H}=A$, the values are real; and if they are all real, $q_{A^{\mathsf H}}=q_A$, and $A^{\mathsf H}=A$ by (ii).
\end{proof}

The real and the complex cases differ in (ii). Over $\mathbb R$, the antisymmetric part of a matrix is invisible in its form, and the rotation $R_{\pi/2}$, with the rows $(0,-1)$ and $(1,0)$, has $\langle\mathbf x,R_{\pi/2}\mathbf x\rangle=0$ for every real $\mathbf x$. Over $\mathbb C$ it is visible: for $\mathbf x=(1,i)$, $R_{\pi/2}\mathbf x=(-i,1)$ and $\mathbf x^{\mathsf H}R_{\pi/2}\mathbf x=1\cdot(-i)+(-i)\cdot1=-2i$, which is not even real. Part (iii) explains why Hermitian matrices are the right class over $\mathbb C$: they are exactly the matrices whose forms take real values, so that it makes sense to ask whether the values are positive.

\subsection{Change of Variables and Congruence}

If the variables are changed by $\mathbf x=S\mathbf y$ with an invertible matrix $S$, the form becomes
\[
q_A(S\mathbf y)=\langle S\mathbf y,AS\mathbf y\rangle=\langle\mathbf y,S^{\mathsf H}AS\mathbf y\rangle=q_{S^{\mathsf H}AS}(\mathbf y),
\]
by \cref{prop:hermitian-inner-product}\,(v): in the new variables the same form has the matrix $S^{\mathsf H}AS$.

\begin{definition}[Congruence]
\label{def:congruence}
Matrices $A,B\in\mathbb F^{n\times n}$ are \term{congruent} if $B=S^{\mathsf H}AS$ for an invertible $S\in\mathbb F^{n\times n}$.
\end{definition}

Congruence is an equivalence relation, since $A=I^{\mathsf H}AI$, $A=(S^{-1})^{\mathsf H}B\,S^{-1}$, and $C=T^{\mathsf H}BT$ gives $C=(ST)^{\mathsf H}A(ST)$; a matrix congruent to a Hermitian matrix is Hermitian, since $(S^{\mathsf H}AS)^{\mathsf H}=S^{\mathsf H}A^{\mathsf H}S$. Congruence must not be confused with similarity, $B=S^{-1}AS$, which describes the same linear map in another basis (\cref{sec:similarity}). Similarity preserves eigenvalues and congruence does not: for $A=I_2$ and $S=\operatorname{diag}(2,1)$, the congruent matrix $S^{\mathsf T}IS=\operatorname{diag}(4,1)$ has the eigenvalues $4$ and $1$, whereas $S^{-1}IS=I$. What congruence preserves is the sign behaviour of the form, since $q_A$ and $q_{S^{\mathsf H}AS}$ take the same values, at the points $S\mathbf y$ and $\mathbf y$; \cref{sec:inertia} makes this precise. The two relations coincide when $S$ is unitary, $S^{\mathsf H}=S^{-1}$, and this is why the spectral theorem of the next section, which finds a unitary $S$ for which $S^{\mathsf H}AS$ is diagonal, simplifies at the same time the linear map $\mathbf x\mapsto A\mathbf x$ and the quadratic form $q_A$.

Completing the square is a congruence. For $a\ne0$, the identity \veqref{1}{eq:complete-square} of \vcref{1}{prop:definite-2x2} says
\[
ah^2+2bhk+ck^2=a\Bigl(h+\frac bak\Bigr)^2+\frac{ac-b^2}a\,k^2 ,
\]
that is, in the variables $y_1=h+bk/a$ and $y_2=k$, the matrix of the form becomes $\operatorname{diag}(a,(ac-b^2)/a)$. The change of variables is $\mathbf y=L^{\mathsf T}\mathbf x$ with the unit lower triangular $L$ whose entry below the diagonal is $b/a$, and the identity is the factorisation $A=LDL^{\mathsf T}$ of \cref{prop:ldu}; \cref{sec:positive-definite} carries this out in $n$ variables.

\subsection{Level Sets in the Plane}

The shape of a quadratic form of two variables is seen in its graph and in its level curves $q(\mathbf x)=c$. \Cref{fig:quadratic-forms} shows three forms,
\begin{equation}
q_1=2x^2+2xy+2y^2,\qquad q_2=x^2+4xy+y^2,\qquad q_3=x^2+2xy+y^2=(x+y)^2 ,
\label{eq:three-forms}
\end{equation}
with the matrices that have the rows $(2,1)$ and $(1,2)$, the rows $(1,2)$ and $(2,1)$, and the rows $(1,1)$ and $(1,1)$. The first is positive for $\mathbf x\ne\mathbf 0$; its graph is a bowl and its level curves are ellipses. The second takes both signs; its graph is a saddle, its positive and negative level curves are hyperbolas, and its zero set is a pair of lines through the origin, the asymptotes of the hyperbolas. The third is nonnegative but vanishes on the line $y=-x$; its graph is a trough and its level curves are pairs of parallel lines. In all three cases the curves are symmetric with respect to the diagonals $y=x$ and $y=-x$, which are the eigenvector lines of the three matrices: the first has the eigenvalues $3$ and $1$, the second $3$ and $-1$, and the third $2$ and $0$, with the eigenvectors $(1,1)$ and $(1,-1)$ in each case. The signs of the eigenvalues decide which of the three pictures one sees, and \cref{sec:principal-axes} proves this in every dimension.

\begin{figure}
\newcommand\sdsurfpanel[1]{\begin{tikzpicture}
\begin{axis}[view={35}{28}, hide axis, width=0.39\textwidth, height=0.36\textwidth,
  xmin=-1, xmax=1, ymin=-1, ymax=1, clip=false]
\addplot3[surf, shader=flat, fill=white, draw=black!45, line width=0.25pt, z buffer=sort, domain=-1:1, y domain=-1:1, samples=11, samples y=11] {#1};
\end{axis}
\end{tikzpicture}}
\newcommand\sdlevelaxis[1]{\begin{tikzpicture}
\begin{axis}[sdplot, width=0.36\textwidth, height=0.36\textwidth, axis equal image,
  xmin=-1.5, xmax=1.5, ymin=-1.5, ymax=1.5, xtick={-1,0,1}, ytick={-1,0,1}, xlabel={$x$}, ylabel={$y$}, clip=true]
#1
\end{axis}
\end{tikzpicture}}
\sdpanel{A}{\sdsurfpanel{2*x^2+2*x*y+2*y^2}}\hfill
\sdpanel{B}{\sdsurfpanel{x^2+4*x*y+y^2}}\hfill
\sdpanel{C}{\sdsurfpanel{(x+y)^2}}
\sdpanelrow
\sdpanel{D}{\sdlevelaxis{%
\foreach \c in {0.25,0.75,1.5,2.5} {
\addplot[black, line width=0.6pt, domain=0:360, samples=90] ({(sqrt(\c/3)*cos(x)+sqrt(\c)*sin(x))/sqrt(2)}, {(sqrt(\c/3)*cos(x)-sqrt(\c)*sin(x))/sqrt(2)});}}}\hfill
\sdpanel{E}{\sdlevelaxis{%
\foreach \c in {0.5,1.5,3} {
\addplot[black, line width=0.6pt, domain=-2.2:2.2, samples=60] ({(sqrt(\c/3)*cosh(x)+sqrt(\c)*sinh(x))/sqrt(2)}, {(sqrt(\c/3)*cosh(x)-sqrt(\c)*sinh(x))/sqrt(2)});
\addplot[black, line width=0.6pt, domain=-2.2:2.2, samples=60] ({(-sqrt(\c/3)*cosh(x)+sqrt(\c)*sinh(x))/sqrt(2)}, {(-sqrt(\c/3)*cosh(x)-sqrt(\c)*sinh(x))/sqrt(2)});
\addplot[black!50, line width=0.6pt, domain=-2.2:2.2, samples=60] ({(sqrt(\c/3)*sinh(x)+sqrt(\c)*cosh(x))/sqrt(2)}, {(sqrt(\c/3)*sinh(x)-sqrt(\c)*cosh(x))/sqrt(2)});
\addplot[black!50, line width=0.6pt, domain=-2.2:2.2, samples=60] ({(sqrt(\c/3)*sinh(x)-sqrt(\c)*cosh(x))/sqrt(2)}, {(sqrt(\c/3)*sinh(x)+sqrt(\c)*cosh(x))/sqrt(2)});}
\addplot[black, dashed, line width=0.5pt, domain=-2:2, samples=2] ({(x+sqrt(3)*x)/sqrt(2)}, {(x-sqrt(3)*x)/sqrt(2)});
\addplot[black, dashed, line width=0.5pt, domain=-2:2, samples=2] ({(x-sqrt(3)*x)/sqrt(2)}, {(x+sqrt(3)*x)/sqrt(2)});}}\hfill
\sdpanel{F}{\sdlevelaxis{%
\foreach \c in {0.25,0.75,1.5} {
\addplot[black, line width=0.6pt, domain=-2:2, samples=2] ({x}, {sqrt(\c)-x});
\addplot[black, line width=0.6pt, domain=-2:2, samples=2] ({x}, {-sqrt(\c)-x});}
\addplot[black, dashed, line width=0.5pt, domain=-2:2, samples=2] ({x}, {-x});}}
\caption{The three quadratic forms \eqref{eq:three-forms} over the square $[-1,1]^2$ (top) and their level curves (bottom). (A, D) The positive definite form $q_1$, with the eigenvalues $3$ and $1$: a bowl, whose level curves $q_1=c$ for $c=0.25,0.75,1.5,2.5$ are ellipses. (B, E) The indefinite form $q_2$, with the eigenvalues $3$ and $-1$: a saddle, whose level curves for $c=0.5,1.5,3$ (black) and $c=-0.5,-1.5,-3$ (grey) are hyperbolas with the zero set (dashed) as asymptotes. (C, F) The positive semidefinite form $q_3=(x+y)^2$, with the eigenvalues $2$ and $0$: a trough, whose level curves for $c=0.25,0.75,1.5$ are pairs of parallel lines and which vanishes on the dashed line.\label{fig:quadratic-forms}}
\end{figure}

\subsection{Where Quadratic Forms Come From}

Quadratic forms arise whenever a quantity depends on a vector to second order. The most important source is Taylor's theorem: by \vcref{1}{thm:taylor-rn}, near a point $\mathbf c$ a $C^2$ function satisfies
\[
f(\mathbf c+\mathbf h)=f(\mathbf c)+\langle\nabla f(\mathbf c),\mathbf h\rangle+\tfrac12\langle\mathbf h,\nabla^2f(\mathbf c)\mathbf h\rangle+o(\|\mathbf h\|^2),
\]
and the Hessian $\nabla^2f(\mathbf c)$ is symmetric because mixed partial derivatives commute (\vcref{1}{sec:partial-derivatives}). At a critical point the quadratic form of the Hessian decides the local behaviour, which is the subject of the second-derivative test (\vcref{1}{thm:second-derivative-test-rn}).

In mechanics, energies are quadratic forms. Consider $n$ masses $m$ on a line, joined to each other and to two fixed walls by $n+1$ equal springs with the constant $k$, and let $x_i$ be the displacement of the $i$th mass from its rest position. The $i$th spring is stretched by $x_{i}-x_{i-1}$, with $x_0=x_{n+1}=0$ for the walls, and the potential energy is
\[
V(\mathbf x)=\frac k2\sum_{i=0}^{n}(x_{i+1}-x_i)^2=\frac k2\Bigl(x_1^2+\sum_{i=1}^{n-1}(x_{i+1}-x_i)^2+x_n^2\Bigr)=\frac k2\langle\mathbf x,T_n\mathbf x\rangle ,
\]
since $\langle\mathbf x,T_n\mathbf x\rangle=2\sum_ix_i^2-2\sum_{i<n}x_ix_{i+1}$ and each squared difference contributes $x_i^2+x_{i+1}^2-2x_ix_{i+1}$. The \term{stiffness matrix} of the chain is $kT_n$, the second difference matrix of \cref{sec:special-matrices} times the spring constant, and the kinetic energy $(m/2)\sum_i\dot x_i^2$ is the quadratic form of the diagonal \term{mass matrix} $mI$; for systems with constraints, such as a double pendulum, the kinetic energy is a quadratic form $\langle\dot{\mathbf x},M\dot{\mathbf x}\rangle/2$ with a symmetric matrix $M$ that need not be diagonal and may depend on the position. \Cref{sec:inertia} returns to the vibrations of such systems.

In statistics, the variance of a linear combination $\langle\mathbf w,\mathbf X\rangle$ of random variables $X_1,\ldots,X_n$ is the quadratic form $\langle\mathbf w,\Sigma\mathbf w\rangle$ of their covariance matrix $\Sigma$ (Volume~IV), and in least squares the function to be minimised is
\[
\|A\mathbf x-\mathbf b\|^2=\langle\mathbf x,A^{\mathsf H}A\mathbf x\rangle-2\operatorname{Re}\langle A^{\mathsf H}\mathbf b,\mathbf x\rangle+\|\mathbf b\|^2 ,
\]
a quadratic form of the Gram matrix $A^{\mathsf H}A$ plus terms of lower degree. In each case the matrix is symmetric or Hermitian, and in each case one asks whether the form is positive, how large and how small it can be on the unit sphere, and what it looks like in well-chosen coordinates.

\section{The Spectral Theorem}
\label{sec:spectral-theorem}

\subsection{Hermitian Matrices Can Be Diagonalised by Unitary Matrices}

The spectral theorem is proved by induction on the dimension: one eigenvector is found, and the orthogonal complement of its line is shown to be invariant. The eigenvector can be found from the fundamental theorem of algebra (\cref{thm:eigenvalues-exist}), but there is a more instructive way, which works over $\mathbb R$ and over $\mathbb C$ alike and which \vcref{1}{sec:lagrange} anticipated with Lagrange multipliers: the point at which the quadratic form is largest on the unit sphere is an eigenvector.

\begin{lemma}[The maximum of a form on the sphere]
\label{lem:variational-eigenvector}
Let $A\in\mathbb F^{n\times n}$ be Hermitian, $n\ge1$. The real function $q_A$ attains a largest and a smallest value on the unit sphere $\{\mathbf x\in\mathbb F^n:\|\mathbf x\|=1\}$, and every unit vector at which one of them is attained is an eigenvector of $A$, with the eigenvalue $q_A(\mathbf x)$. In particular, $A$ has an eigenvector in $\mathbb F^n$, for a real eigenvalue.
\end{lemma}

\begin{proof}
Writing each complex coordinate by its real and imaginary parts identifies $\mathbb C^n$ with $\mathbb R^{2n}$, with the same length; the unit sphere is then closed and bounded, and $q_A$ is a polynomial in the real coordinates, real-valued by \cref{prop:form-matrix}\,(iii). By the extreme value theorem (\vcref{1}{thm:evt-rn}) it attains a largest value at some unit vector $\mathbf x$. Let $\mathbf y\perp\mathbf x$ and $t\in\mathbb R$. Then $\|\mathbf x+t\mathbf y\|^2=1+t^2\|\mathbf y\|^2$ by Pythagoras' theorem, and as in the proof of \cref{prop:form-matrix},
\[
q_A(\mathbf x+t\mathbf y)=q_A(\mathbf x)+2t\operatorname{Re}\langle\mathbf x,A\mathbf y\rangle+t^2q_A(\mathbf y),
\]
since $\mathbf y^{\mathsf H}A\mathbf x=\overline{\mathbf x^{\mathsf H}A^{\mathsf H}\mathbf y}=\overline{\mathbf x^{\mathsf H}A\mathbf y}$. The function
\[
\varphi(t)\coloneqq q_A\Bigl(\frac{\mathbf x+t\mathbf y}{\|\mathbf x+t\mathbf y\|}\Bigr)=\frac{q_A(\mathbf x)+2t\operatorname{Re}\langle\mathbf x,A\mathbf y\rangle+t^2q_A(\mathbf y)}{1+t^2\|\mathbf y\|^2}
\]
has a maximum at $t=0$, and it is differentiable, so $\varphi'(0)=2\operatorname{Re}\langle\mathbf x,A\mathbf y\rangle=0$ by Fermat's theorem (\vcref{1}{thm:fermat-extremum}). If $\mathbb F=\mathbb C$, the same holds for $i\mathbf y$, which is also orthogonal to $\mathbf x$, and $\operatorname{Re}\langle\mathbf x,A(i\mathbf y)\rangle=-\operatorname{Im}\langle\mathbf x,A\mathbf y\rangle$. Hence $\langle A\mathbf x,\mathbf y\rangle=\langle\mathbf x,A\mathbf y\rangle=0$ for every $\mathbf y\perp\mathbf x$, that is, $A\mathbf x$ lies in the orthogonal complement of $\{\mathbf x\}^\perp$, which is the line spanned by $\mathbf x$ (\cref{thm:orthogonal-decomposition}). So $A\mathbf x=\lambda\mathbf x$ for some $\lambda$, and $\lambda=\langle\mathbf x,\lambda\mathbf x\rangle=q_A(\mathbf x)$ is real. For the smallest value apply the argument to $-A$.
\end{proof}

The proof is the method of Lagrange multipliers for this particular problem, carried out by hand: the derivative of $q_A$ along every direction tangent to the sphere vanishes at the maximum, which says that $\nabla q_A(\mathbf x)=2A\mathbf x$ is a multiple of the normal vector $\mathbf x$. It needs neither the implicit function theorem nor the fundamental theorem of algebra.

\begin{theorem}[Spectral theorem]
\label{thm:spectral}
Let $A\in\mathbb C^{n\times n}$ be Hermitian. There are a unitary matrix $Q\in\mathbb C^{n\times n}$ and a real diagonal matrix $\Lambda=\operatorname{diag}(\lambda_1,\ldots,\lambda_n)$ with $\lambda_1\le\cdots\le\lambda_n$ such that
\[
A=Q\Lambda Q^{\mathsf H}.
\]
The numbers $\lambda_j$ are the eigenvalues of $A$, each repeated according to its algebraic multiplicity, and the columns $\mathbf q_1,\ldots,\mathbf q_n$ of $Q$ form an orthonormal basis of $\mathbb C^n$ with $A\mathbf q_j=\lambda_j\mathbf q_j$. If $A$ is real and symmetric, $Q$ can be chosen real, so that $A=Q\Lambda Q^{\mathsf T}$ with a real orthogonal $Q$. Conversely, $Q\Lambda Q^{\mathsf H}$ is Hermitian for every unitary $Q$ and every real diagonal $\Lambda$, and $Q\Lambda Q^{\mathsf T}$ is real and symmetric for every real orthogonal $Q$ and real diagonal $\Lambda$.
\end{theorem}

\begin{proof}
We prove by induction on $n$ that there are a unitary $Q$ and a real diagonal $\Lambda$ with $A=Q\Lambda Q^{\mathsf H}$, real when $A$ is real. For $n=1$, $A=[a_{11}]$ with $a_{11}$ real, and $Q=[1]$. Let $n\ge2$. By \cref{lem:variational-eigenvector}, applied with $\mathbb F=\mathbb C$, or with $\mathbb F=\mathbb R$ if $A$ is real, there is a unit vector $\mathbf x$ with $A\mathbf x=\lambda\mathbf x$, $\lambda$ real. Extend $\mathbf x$ to an orthonormal basis $\mathbf x,\mathbf w_2,\ldots,\mathbf w_n$ of $\mathbb F^n$ (\cref{cor:orthonormal-basis}), and let $W$ be the $n\times(n-1)$ matrix with the columns $\mathbf w_j$, so that $U\coloneqq[\,\mathbf x\ W\,]$ is unitary. By block multiplication,
\[
U^{\mathsf H}AU=\begin{bmatrix}\mathbf x^{\mathsf H}A\mathbf x&\mathbf x^{\mathsf H}AW\\W^{\mathsf H}A\mathbf x&W^{\mathsf H}AW\end{bmatrix}=\begin{bmatrix}\lambda&\mathbf 0^{\mathsf T}\\\mathbf 0&B\end{bmatrix},\qquad B\coloneqq W^{\mathsf H}AW ,
\]
because $W^{\mathsf H}A\mathbf x=\lambda W^{\mathsf H}\mathbf x=\mathbf 0$, the columns of $W$ being orthogonal to $\mathbf x$, and $\mathbf x^{\mathsf H}AW=(A\mathbf x)^{\mathsf H}W=\lambda\mathbf x^{\mathsf H}W=\mathbf 0^{\mathsf T}$, since $A^{\mathsf H}=A$ and $\lambda$ is real. The matrix $B$ is Hermitian of order $n-1$, real if $A$ is, and by the induction hypothesis $B=Q_2\Lambda_2Q_2^{\mathsf H}$ with a unitary $Q_2$ and a real diagonal $\Lambda_2$. Then
\[
A=U\begin{bmatrix}1&\mathbf 0^{\mathsf T}\\\mathbf 0&Q_2\end{bmatrix}\begin{bmatrix}\lambda&\mathbf 0^{\mathsf T}\\\mathbf 0&\Lambda_2\end{bmatrix}\begin{bmatrix}1&\mathbf 0^{\mathsf T}\\\mathbf 0&Q_2\end{bmatrix}^{\mathsf H}U^{\mathsf H},
\]
and the product of $U$ and the block diagonal matrix is unitary, as a product of unitary matrices (\cref{prop:orthonormal-columns}\,(iv)). This completes the induction.

A permutation matrix $P$ is real orthogonal, and $P^{\mathsf T}\Lambda P$ is diagonal with the diagonal entries of $\Lambda$ rearranged (\cref{prop:permutation-matrices}); replacing $Q$ by $QP$ and $\Lambda$ by $P^{\mathsf T}\Lambda P$ puts the diagonal in increasing order. Since $Q^{\mathsf H}=Q^{-1}$, $A$ is similar to $\Lambda$, so $p_A=p_\Lambda=\prod_j(x-\lambda_j)$ (\cref{prop:charpoly-similar}), and the $\lambda_j$ are the eigenvalues with their algebraic multiplicities; $AQ=Q\Lambda$ says $A\mathbf q_j=\lambda_j\mathbf q_j$, and the columns of a unitary matrix are orthonormal. Conversely, $(Q\Lambda Q^{\mathsf H})^{\mathsf H}=Q\Lambda^{\mathsf H}Q^{\mathsf H}=Q\Lambda Q^{\mathsf H}$ because $\Lambda$ is real and diagonal.
\end{proof}

The theorem settles the question that \cref{ch:eigenvalues} left open: every Hermitian matrix is diagonalisable, and by a unitary matrix. In particular the geometric multiplicity of each eigenvalue equals its algebraic multiplicity (\cref{thm:diagonalisable}), and eigenvalues of Hermitian matrices are never defective. Together with \cref{prop:hermitian-eigen} the theorem gives the complete picture: the eigenspaces are orthogonal to each other, their dimensions add up to $n$, and an orthonormal basis of each, chosen by the Gram--Schmidt process if necessary, gives the columns of $Q$. The hypotheses of the convergence theorem of the power method (\cref{thm:power-method}) and of the Bauer--Fike theorem (\cref{thm:bauer-fike}), which assumed diagonalisability, therefore hold for every Hermitian matrix, and the Bauer--Fike bound holds with $\kappa_2(Q)=1$.

The converse part shows that real symmetry is exactly what is needed for an orthonormal basis of real eigenvectors: if a real matrix has such a basis, it is $Q\Lambda Q^{\mathsf T}$ with a real orthogonal $Q$ and a real diagonal $\Lambda$, hence symmetric. The theorem is equally a statement about operators.

\begin{corollary}[Self-adjoint operators]
\label{cor:spectral-operator}
Let $T$ be a self-adjoint operator on a finite-dimensional inner product space $V$ over $\mathbb F$. Then the eigenvalues of $T$ are real, and $V$ has an orthonormal basis consisting of eigenvectors of $T$. Conversely, an operator with an orthonormal basis of eigenvectors whose eigenvalues are real is self-adjoint.
\end{corollary}

\begin{proof}
Let $\mathcal B$ be an orthonormal basis of $V$ (\cref{cor:orthonormal-basis}). By \cref{prop:adjoint-matrix}, $A\coloneqq[T]_{\mathcal B}$ is Hermitian, so $A=Q\Lambda Q^{\mathsf H}$ by \cref{thm:spectral}. Let $\mathbf v_j$ be the vector with the coordinate vector $\mathbf q_j$. Since $\mathcal B$ is orthonormal, $\langle\mathbf v_i,\mathbf v_j\rangle=\langle\mathbf q_i,\mathbf q_j\rangle=\delta_{ij}$ (\cref{thm:orthonormal-expansion}), and $\mathbf v_j$ is an eigenvector of $T$ for $\lambda_j$ by \cref{prop:similar-eigen}. Conversely, the matrix of such an operator in its orthonormal basis of eigenvectors is real and diagonal, hence Hermitian, and the operator is self-adjoint by \cref{prop:adjoint-matrix}.
\end{proof}

The second derivative with a sign, $-D^2$, on the span of $\sin x,\sin2x,\ldots,\sin Nx$ in $C[0,\pi]$ with the inner product $\int_0^\pi fg$, is self-adjoint by \cref{sec:adjoints}, and the functions $\sin kx$ are eigenvectors for the eigenvalues $k^2$; they are orthogonal by \cref{prop:trig-orthogonal}, and divided by their norm $\sqrt{\pi/2}$ they form the orthonormal basis of the corollary. In infinite dimensions the spectral theorem needs more care, as the note below explains.

\begin{note}{The spectral theorem in infinite dimensions}
On an infinite-dimensional inner product space a self-adjoint operator need not have any eigenvector. Multiplication by $x$ on $C[0,1]$ with the inner product $\int_0^1fg$ is self-adjoint, since $\int_0^1(xf)g=\int_0^1f(xg)$, but $xf(x)=\lambda f(x)$ for all $x$ forces $f=0$ away from $x=\lambda$, hence everywhere by continuity. The proof of \cref{lem:variational-eigenvector} fails because the unit sphere is no longer compact, and the maximum of the form need not be attained. The theorem survives for \emph{compact} self-adjoint operators on Hilbert spaces, those that map bounded sets into sets with compact closure, such as integral operators $f\mapsto\int k(x,t)f(t)\,dt$ with a continuous symmetric kernel $k$: the closure of their image has an orthonormal basis consisting of eigenvectors, for real eigenvalues that tend to zero if there are infinitely many of them~\cite{kreyszig1978}. This is the theory that Hilbert developed from integral equations~\cite{hilbert1904}. The inverses of the second derivative and of more general Sturm--Liouville operators $-(pu')'+ru$ with boundary conditions, where they exist, are integral operators of this kind, which is why such differential operators have orthogonal bases of eigenfunctions, the sines of the example above among them. Volume~III develops this theory.
\end{note}

\subsection{The Spectral Decomposition}

Reading the product $Q\Lambda Q^{\mathsf H}$ by columns and rows (\cref{thm:four-readings}) writes a Hermitian matrix as a sum of rank-one matrices, and grouping the terms by eigenvalues gives the spectral projections of \cref{lem:spectral-projections}, which are now orthogonal projections.

\begin{proposition}[Spectral decomposition]
\label{prop:spectral-decomposition}
Let $A=Q\Lambda Q^{\mathsf H}$ be Hermitian as in \cref{thm:spectral}, let $\mu_1,\ldots,\mu_k$ be its pairwise different eigenvalues, and let $P_i$ be the orthogonal projection onto the eigenspace $E_{\mu_i}$. Then
\[
A=\sum_{j=1}^n\lambda_j\,\mathbf q_j\mathbf q_j^{\mathsf H}=\sum_{i=1}^k\mu_iP_i,\qquad P_i=\sum_{j:\,\lambda_j=\mu_i}\mathbf q_j\mathbf q_j^{\mathsf H},
\]
with $P_1+\cdots+P_k=I$, $P_iP_l=0$ for $i\ne l$, and $p(A)=\sum_ip(\mu_i)P_i$ for every polynomial $p$. The eigenspaces are mutually orthogonal, $\mathbb F^n=E_{\mu_1}\oplus\cdots\oplus E_{\mu_k}$, and the projections $P_i$ are the spectral projections $\ell_i(A)$ of \cref{lem:spectral-projections}; they do not depend on the choice of $Q$.
\end{proposition}

\begin{proof}
The first sum is the product $Q\Lambda Q^{\mathsf H}$ read as the sum of the outer products of the columns $\lambda_j\mathbf q_j$ of $Q\Lambda$ with the rows $\mathbf q_j^{\mathsf H}$ of $Q^{\mathsf H}$. The $\operatorname{am}(\mu_i)$ vectors $\mathbf q_j$ with $\lambda_j=\mu_i$ are orthonormal vectors of $E_{\mu_i}$, so $\dim E_{\mu_i}\ge\operatorname{am}(\mu_i)$, and by \cref{thm:gm-am} they form an orthonormal basis of $E_{\mu_i}$. The sum of their outer products is therefore the orthogonal projection onto $E_{\mu_i}$ by \cref{thm:projection-matrix}\,(ii), which gives the second form and shows that $P_i$ depends only on $E_{\mu_i}$. The eigenspaces are orthogonal by \cref{prop:hermitian-eigen}\,(ii), and the matrix is diagonalisable, so they form a direct sum equal to $\mathbb F^n$ (\cref{thm:diagonalisable}). By \cref{lem:spectral-projections}\,(i), $\ell_i(A)$ maps $\mathbf v=\mathbf v_1+\cdots+\mathbf v_k$ with $\mathbf v_l\in E_{\mu_l}$ to $\mathbf v_i$, and since the summands are orthogonal, $\mathbf v_i$ is the orthogonal projection of $\mathbf v$ onto $E_{\mu_i}$. The remaining identities are \cref{lem:spectral-projections}\,(ii).
\end{proof}

The decomposition says that $A$ acts on the orthogonal pieces $E_{\mu_i}$ of $\mathbb F^n$ separately, multiplying the piece of $\mathbf v$ in $E_{\mu_i}$ by $\mu_i$. For a general diagonalisable matrix the spectral projections are oblique, as for the matrix $B$ of \eqref{eq:two-by-two-examples}, whose projection $P_1$ maps the plane onto the $x_1$-axis along the line spanned by $(1,1)$; for Hermitian matrices they are orthogonal, and this is what makes Hermitian matrices so much easier to work with. The formula $p(A)=\sum_ip(\mu_i)P_i$ suggests a definition for functions other than polynomials: for a function $f$ defined at the eigenvalues, $f(A)\coloneqq\sum_if(\mu_i)P_i=Qf(\Lambda)Q^{\mathsf H}$, with $f(\Lambda)=\operatorname{diag}(f(\lambda_1),\ldots,f(\lambda_n))$. The square root of \cref{sec:cholesky} is the first example, and \cref{ch:canonical-forms} defines functions of general matrices, among them the exponential $e^A$.

As an example, let $\mathbf w=(1,1,1)$ and
\[
A=I+\mathbf w\mathbf w^{\mathsf T}=\begin{bmatrix}2&1&1\\1&2&1\\1&1&2\end{bmatrix}.
\]
Since $A\mathbf w=\mathbf w+3\mathbf w=4\mathbf w$ and $A\mathbf x=\mathbf x$ for every $\mathbf x\perp\mathbf w$, the eigenvalues are $1$, $1$ and $4$, the eigenspace $E_4$ is the line spanned by $\mathbf w$, and $E_1$ is the plane $x_1+x_2+x_3=0$. The vectors $(1,-1,0)$ and $(1,0,-1)$ form a basis of $E_1$, but not an orthogonal one; the Gram--Schmidt process turns them into $\mathbf q_1=(1,-1,0)/\sqrt2$ and, since $(1,0,-1)-\langle\mathbf q_1,(1,0,-1)\rangle\mathbf q_1=(1/2,1/2,-1)$, into $\mathbf q_2=(1,1,-2)/\sqrt6$. With $\mathbf q_3=\mathbf w/\sqrt3$,
\[
Q=\begin{bmatrix}1/\sqrt2&1/\sqrt6&1/\sqrt3\\-1/\sqrt2&1/\sqrt6&1/\sqrt3\\0&-2/\sqrt6&1/\sqrt3\end{bmatrix},\qquad\Lambda=\begin{bmatrix}1&0&0\\0&1&0\\0&0&4\end{bmatrix},\qquad A=Q\Lambda Q^{\mathsf T},
\]
and the spectral decomposition is $A=1\cdot P_1+4\cdot P_2$, with the projection $P_2=\mathbf w\mathbf w^{\mathsf T}/3$ onto $E_4$ and the projection $P_1=I-\mathbf w\mathbf w^{\mathsf T}/3$ onto $E_1$. Any other orthonormal basis of the plane $E_1$ would serve for $\mathbf q_1$ and $\mathbf q_2$, so $Q$ is not unique, but $P_1$ is the same for all of them.

\subsection{Normal Matrices}

The proof of \cref{thm:spectral} used the Hermitian property only to show that the orthogonal complement of an eigenvector is invariant. A weaker hypothesis suffices for this, and it characterises the matrices that can be diagonalised by unitary matrices. A matrix $A\in\mathbb C^{n\times n}$ is \term{normal} if $A^{\mathsf H}A=AA^{\mathsf H}$ (\cref{def:operator-classes}).

\begin{theorem}[Unitary diagonalisation]
\label{thm:normal-unitary}
A matrix $A\in\mathbb C^{n\times n}$ can be written as $A=Q\Lambda Q^{\mathsf H}$ with a unitary $Q$ and a diagonal $\Lambda\in\mathbb C^{n\times n}$ if and only if $A$ is normal. A normal matrix is Hermitian if and only if its eigenvalues are real, unitary if and only if they have modulus $1$, and skew-Hermitian if and only if they lie on the imaginary axis.
\end{theorem}

\begin{proof}
If $A=Q\Lambda Q^{\mathsf H}$, then $A^{\mathsf H}=Q\bar\Lambda Q^{\mathsf H}$, and $A^{\mathsf H}A=Q\bar\Lambda\Lambda Q^{\mathsf H}=Q\Lambda\bar\Lambda Q^{\mathsf H}=AA^{\mathsf H}$. For the converse we use induction on $n$, as in \cref{thm:spectral}. Let $A$ be normal, $n\ge2$. By \cref{thm:eigenvalues-exist}, $A$ has an eigenvalue $\lambda$ and a unit eigenvector $\mathbf x$. The matrix $N\coloneqq A-\lambda I$ is normal too, since $N^{\mathsf H}N$ and $NN^{\mathsf H}$ both equal $A^{\mathsf H}A-\bar\lambda A-\lambda A^{\mathsf H}+|\lambda|^2I$, so $\|N^{\mathsf H}\mathbf x\|^2=\mathbf x^{\mathsf H}NN^{\mathsf H}\mathbf x=\mathbf x^{\mathsf H}N^{\mathsf H}N\mathbf x=\|N\mathbf x\|^2=0$, that is, $A^{\mathsf H}\mathbf x=\bar\lambda\mathbf x$. With $U=[\,\mathbf x\ W\,]$ unitary as before, $W^{\mathsf H}A\mathbf x=\lambda W^{\mathsf H}\mathbf x=\mathbf 0$ and $\mathbf x^{\mathsf H}AW=(A^{\mathsf H}\mathbf x)^{\mathsf H}W=\lambda\mathbf x^{\mathsf H}W=\mathbf 0^{\mathsf T}$, so $U^{\mathsf H}AU=\operatorname{diag}(\lambda,B)$ with $B=W^{\mathsf H}AW$. A matrix unitarily similar to a normal matrix is normal, by the first computation of this proof with $\Lambda$ replaced by any matrix, and comparing the lower right blocks of $\operatorname{diag}(\lambda,B)^{\mathsf H}\operatorname{diag}(\lambda,B)$ and $\operatorname{diag}(\lambda,B)\operatorname{diag}(\lambda,B)^{\mathsf H}$ gives $B^{\mathsf H}B=BB^{\mathsf H}$. The induction hypothesis applied to $B$ completes the proof as before. Finally, for $A=Q\Lambda Q^{\mathsf H}$ the conditions $A^{\mathsf H}=A$, $A^{\mathsf H}A=I$ and $A^{\mathsf H}=-A$ are equivalent to $\bar\Lambda=\Lambda$, $\bar\Lambda\Lambda=I$ and $\bar\Lambda=-\Lambda$.
\end{proof}

The rotation $R_\varphi$ is normal, being orthogonal, and it is diagonalised by the unitary matrix with the columns $(1,-i)/\sqrt2$ and $(1,i)/\sqrt2$, the normalised eigenvectors of \cref{sec:diagonalisation}; over $\mathbb R$ it cannot be diagonalised at all, because its eigenvalues are not real. For real matrices, therefore, the real form of the theorem holds only for symmetric matrices. The matrices $B$ of \eqref{eq:two-by-two-examples} and $A_2$ of \cref{sec:matrix-powers} are not normal, the first being diagonalisable but not by a unitary matrix, and the second not diagonalisable at all. \Cref{ch:canonical-forms} obtains the theorem again from the Schur form and describes the real normal matrices by a real block diagonal form.

\subsection{Commuting Hermitian Matrices}

\Cref{thm:simultaneous} diagonalised two commuting diagonalisable matrices with the same invertible matrix. For Hermitian matrices the common matrix can be chosen unitary.

\begin{theorem}[Simultaneous unitary diagonalisation]
\label{thm:simultaneous-unitary}
Let $A,B\in\mathbb F^{n\times n}$ be Hermitian. Then $AB=BA$ if and only if there is a unitary matrix $Q\in\mathbb F^{n\times n}$, real orthogonal if $\mathbb F=\mathbb R$, for which $Q^{\mathsf H}AQ$ and $Q^{\mathsf H}BQ$ are both diagonal, that is, if and only if $\mathbb F^n$ has an orthonormal basis of common eigenvectors of $A$ and $B$.
\end{theorem}

\begin{proof}
Diagonal matrices commute, which gives one direction as in \cref{thm:simultaneous}. Conversely, let $AB=BA$, and let $\mu$ be an eigenvalue of $A$. As in that proof, $E_\mu(A)$ is invariant under $B$, and the restriction of $B$ to $E_\mu(A)$, with the inner product of $\mathbb F^n$, is self-adjoint, since $\langle\mathbf x,B\mathbf y\rangle=\langle B\mathbf x,\mathbf y\rangle$ for all $\mathbf x,\mathbf y$. By \cref{cor:spectral-operator}, $E_\mu(A)$ has an orthonormal basis of eigenvectors of $B$, which are eigenvectors of $A$ as well. The eigenspaces of $A$ are mutually orthogonal and their direct sum is $\mathbb F^n$ (\cref{prop:spectral-decomposition}), so these bases together form an orthonormal basis of $\mathbb F^n$, the columns of $Q$.
\end{proof}

The proof shows how to find $Q$: diagonalise $A$, and inside each eigenspace of $A$ diagonalise $B$. If all eigenvalues of $A$ are simple, each eigenspace is a line, and the eigenvectors of $A$ are automatically eigenvectors of every Hermitian matrix that commutes with $A$.

\subsection{A Picture of the Theorem}

Read from right to left, $A\mathbf x=Q\Lambda Q^{\mathsf H}\mathbf x$ applies three maps in turn. First $Q^{\mathsf H}$ computes the coordinates $y_j=\langle\mathbf q_j,\mathbf x\rangle$ of $\mathbf x$ in the orthonormal basis of eigenvectors; for a real $2\times2$ matrix this is a rotation or a reflection, which preserves lengths and angles and takes $\mathbf q_1$ and $\mathbf q_2$ to $\mathbf e_1$ and $\mathbf e_2$. Then $\Lambda$ stretches each coordinate by its eigenvalue, and finally $Q$ turns the result back. \Cref{fig:spectral-action} follows the unit circle through these steps for
\[
A=R_{\pi/6}\begin{bmatrix}2&0\\0&1/2\end{bmatrix}R_{\pi/6}^{\mathsf T}=\frac18\begin{bmatrix}13&3\sqrt3\\3\sqrt3&7\end{bmatrix},
\]
whose eigenvectors $\mathbf q_1=(\cos(\pi/6),\sin(\pi/6))$ and $\mathbf q_2=(-\sin(\pi/6),\cos(\pi/6))$ have the eigenvalues $2$ and $1/2$. The image of the unit circle is an ellipse whose axes are the eigenvector lines and whose semi-axes are the moduli of the eigenvalues. For the nonsymmetric matrix of \cref{fig:unit-circle-image}B the image is again an ellipse, but its axes are not eigenvector lines; \cref{ch:svd} finds the axes of the image ellipse for every matrix, through the spectral theorem applied to $A^{\mathsf H}A$.

\begin{figure}
\newcommand\sdactionpanel[5]{\begin{tikzpicture}[scale=1.05, font=\scriptsize]
\draw[black!35, line width=0.4pt] (0,0) circle (1);
\draw[line width=0.7pt, smooth, samples=100, domain=0:360] plot ({#1*cos(\x)+#2*sin(\x)}, {#3*cos(\x)+#4*sin(\x)});
\draw[line width=0.4pt, -{Latex[length=1.3mm]}] (-2.25,0) -- (2.35,0);
\draw[line width=0.4pt, -{Latex[length=1.3mm]}] (0,-1.35) -- (0,1.45);
#5
\end{tikzpicture}}
% q1 at 30 degrees, q2 at 120 degrees; the point x at 75 degrees
\sdpanel{A}{\sdactionpanel{1}{0}{0}{1}{%
\draw[line width=0.9pt, -{Latex[length=1.6mm]}] (0,0) -- (0.866,0.5);
\node[below, inner sep=2pt] at (0.58,0.33) {$\mathbf q_1$};
\draw[line width=0.9pt, -{Latex[length=1.6mm]}] (0,0) -- (-0.5,0.866);
\node[left, inner sep=2pt] at (-0.3,0.55) {$\mathbf q_2$};
\fill (0.2588,0.9659) circle (1.5pt);
\node[above right, inner sep=1pt] at (0.27,0.98) {$\mathbf x$};}}\hfill
\sdpanel{B}{\sdactionpanel{1}{0}{0}{1}{%
\draw[line width=0.9pt, -{Latex[length=1.6mm]}] (0,0) -- (1,0);
\node[below, inner sep=2pt] at (0.8,0) {$\mathbf e_1$};
\draw[line width=0.9pt, -{Latex[length=1.6mm]}] (0,0) -- (0,1);
\node[left, inner sep=2pt] at (0,0.8) {$\mathbf e_2$};
\fill (0.7071,0.7071) circle (1.5pt);
\node[above right, inner sep=1pt] at (0.72,0.72) {$Q^{\mathsf T}\mathbf x$};}}
\sdpanelrow
\sdpanel{C}{\sdactionpanel{2}{0}{0}{0.5}{%
\draw[line width=0.9pt, -{Latex[length=1.6mm]}] (0,0) -- (2,0);
\node[below, inner sep=2pt] at (1.45,0) {$2\mathbf e_1$};
\draw[line width=0.9pt, -{Latex[length=1.6mm]}] (0,0) -- (0,0.5);
\node[left, inner sep=2pt] at (0,0.27) {$\mathbf e_2/2$};
\fill (1.4142,0.3536) circle (1.5pt);
\node[above right, inner sep=1pt] at (1.42,0.38) {$\Lambda Q^{\mathsf T}\mathbf x$};}}\hfill
\sdpanel{D}{\sdactionpanel{1.7321}{-0.25}{1}{0.4330}{%
\draw[black, dashed, line width=0.4pt] (-2.25,-1.3) -- (2.25,1.3);
\draw[black, dashed, line width=0.4pt] (0.7,-1.21) -- (-0.7,1.21);
\draw[line width=0.9pt, -{Latex[length=1.6mm]}] (0,0) -- (1.7321,1);
\node[anchor=west, inner sep=2pt] at (1.9,0.78) {$2\mathbf q_1$};
\draw[line width=0.9pt, -{Latex[length=1.6mm]}] (0,0) -- (-0.25,0.433);
\node[anchor=south east, inner sep=1pt] at (-0.74,0.93) {$\mathbf q_2/2$};
\fill (1.0479,1.0133) circle (1.5pt);
\node[above, inner sep=2pt] at (1.0479,1.03) {$A\mathbf x$};}}
\caption{The action of $A=Q\Lambda Q^{\mathsf T}$ with $Q=R_{\pi/6}$ and $\Lambda=\operatorname{diag}(2,1/2)$, followed on the unit circle and on one point $\mathbf x$ (dot); the grey circle is the unit circle, and the black curve is its image after each step. (A) The orthonormal eigenvectors $\mathbf q_1$ and $\mathbf q_2$. (B) After $Q^{\mathsf T}$, a rotation by $-\pi/6$, they have become $\mathbf e_1$ and $\mathbf e_2$, and the circle is unchanged. (C) $\Lambda$ stretches the first coordinate by $2$ and the second by $1/2$, and the circle becomes an ellipse with axes along the coordinate axes. (D) $Q$ rotates back: the image of the circle under $A$ is the ellipse whose axes are the eigenvector lines (dashed), with the semi-axes $2$ and $1/2$, and $A\mathbf x$ is the image of the point.\label{fig:spectral-action}}
\end{figure}

\section{Principal Axes}
\label{sec:principal-axes}

\subsection{Sums of Squares}

In the coordinates of an orthonormal basis of eigenvectors, a Hermitian form has no cross terms.

\begin{theorem}[Principal axes]
\label{thm:principal-axes}
Let $A=Q\Lambda Q^{\mathsf H}$ be Hermitian as in \cref{thm:spectral}, and for $\mathbf x\in\mathbb F^n$ let $\mathbf y\coloneqq Q^{\mathsf H}\mathbf x$, so that $y_j=\langle\mathbf q_j,\mathbf x\rangle$ and $\|\mathbf y\|=\|\mathbf x\|$. Then
\[
q_A(\mathbf x)=\sum_{j=1}^n\lambda_j|y_j|^2,\qquad\text{and hence}\qquad\lambda_1\|\mathbf x\|^2\le q_A(\mathbf x)\le\lambda_n\|\mathbf x\|^2 .
\]
Consequently $A$ is positive definite if and only if all its eigenvalues are positive, positive semidefinite if and only if they are all nonnegative, negative definite or semidefinite if and only if they are all negative or all nonpositive, and indefinite if and only if it has a positive and a negative eigenvalue.
\end{theorem}

\begin{proof}
$q_A(\mathbf x)=\mathbf x^{\mathsf H}Q\Lambda Q^{\mathsf H}\mathbf x=\mathbf y^{\mathsf H}\Lambda\mathbf y=\sum_j\lambda_j\bar y_jy_j$, and $\|Q^{\mathsf H}\mathbf x\|=\|\mathbf x\|$ because $Q^{\mathsf H}$ is unitary. The bounds follow from $\lambda_1\le\lambda_j\le\lambda_n$ and $\sum_j|y_j|^2=\|\mathbf x\|^2$. Since $\mathbf x\mapsto\mathbf y$ is a bijection of $\mathbb F^n$ that maps $\mathbf 0$ to $\mathbf 0$, $q_A$ is positive away from $\mathbf 0$ if and only if $\sum_j\lambda_j|y_j|^2$ is, and this holds if and only if every $\lambda_j>0$, by taking $\mathbf y=\mathbf e_j$ for one direction and adding positive terms for the other. The other cases are the same.
\end{proof}

For a real symmetric $A$ this is \vcref{1}{thm:definiteness-criteria}\,(i), whose proof is now complete: the eigenvalues are real by \cref{prop:hermitian-eigen}, and the sign behaviour of the form is that of its eigenvalues. The lines spanned by the eigenvectors are the \term{principal axes} of the form, a name from mechanics, where the principal axes of a rigid body are the eigenvectors of its symmetric inertia matrix. The theorem is the $n$-dimensional version of completing the square, with an orthogonal change of variables instead of a triangular one.

\subsection{Conics and Quadrics}

The theorem classifies the level sets of real quadratic forms. For a real symmetric $2\times2$ matrix with the eigenvalues $\lambda_1\le\lambda_2$ and the orthonormal eigenvectors $\mathbf q_1,\mathbf q_2$, the curve $\langle\mathbf x,A\mathbf x\rangle=1$ is, in the coordinates $y_1,y_2$ along $\mathbf q_1,\mathbf q_2$, the curve $\lambda_1y_1^2+\lambda_2y_2^2=1$, which is
\begin{itemize}
\item an ellipse with the semi-axes $1/\sqrt{\lambda_1}$ and $1/\sqrt{\lambda_2}$ along $\mathbf q_1$ and $\mathbf q_2$, if $0<\lambda_1\le\lambda_2$, and a circle if $\lambda_1=\lambda_2$;
\item a hyperbola with its vertices on the $\mathbf q_2$-axis at the distance $1/\sqrt{\lambda_2}$ from the origin and the asymptotes $\sqrt{-\lambda_1}\,y_1=\pm\sqrt{\lambda_2}\,y_2$, if $\lambda_1<0<\lambda_2$;
\item a pair of parallel lines $y_2=\pm1/\sqrt{\lambda_2}$, if $\lambda_1=0<\lambda_2$;
\item empty, if $\lambda_2\le0$.
\end{itemize}
The longer axis of an ellipse belongs to the \emph{smaller} eigenvalue: a direction in which the form grows slowly must be followed further to reach the value $1$. For example, the matrix with the rows $(5,2)$ and $(2,8)$ has the trace $13$ and the determinant $36$, hence the eigenvalues $4$ and $9$ by \eqref{eq:eigen-2x2}, with the eigenvectors $(2,-1)$ and $(1,2)$. The ellipse
\[
5x^2+4xy+8y^2=36
\]
of \cref{fig:ellipse-axes} becomes $4y_1^2+9y_2^2=36$, that is $y_1^2/9+y_2^2/4=1$, in the coordinates along $\mathbf q_1=(2,-1)/\sqrt5$ and $\mathbf q_2=(1,2)/\sqrt5$: its semi-axes are $3$ along $\mathbf q_1$ and $2$ along $\mathbf q_2$. In three variables the surfaces $\langle\mathbf x,A\mathbf x\rangle=1$ are, according to the signs of the three eigenvalues, ellipsoids ($+,+,+$), hyperboloids of one sheet ($-,+,+$) or of two sheets ($-,-,+$), cylinders over an ellipse or a hyperbola when an eigenvalue vanishes, or pairs of planes.

\begin{figure}
\sdfloat{%
\begin{tikzpicture}
\begin{axis}[sddiagram, width=0.56\textwidth, height=0.56\textwidth, axis equal image,
  xmin=-3.6, xmax=3.8, ymin=-3.0, ymax=3.1, xtick=\empty, ytick=\empty,
  xlabel={$x$}, ylabel={$y$}]
\addplot[black, dashed, line width=0.4pt, domain=-3.5:3.5, samples=2] ({x}, {-x/2});
\addplot[black, dashed, line width=0.4pt, domain=-1.45:1.45, samples=2] ({x}, {2*x});
\addplot[black, line width=0.8pt, domain=0:360, samples=120] ({(6*cos(x)+2*sin(x))/sqrt(5)}, {(-3*cos(x)+4*sin(x))/sqrt(5)});
\draw[line width=0.9pt, -{Latex[length=1.8mm]}] (axis cs:0,0) -- (axis cs:2.6833,-1.3416);
\draw[line width=0.9pt, -{Latex[length=1.8mm]}] (axis cs:0,0) -- (axis cs:0.8944,1.7889);
\node[font=\scriptsize, below left, inner sep=1pt] at (axis cs:2.6833,-1.42) {$3\mathbf q_1$};
\node[font=\scriptsize, left, inner sep=2pt] at (axis cs:0.85,1.85) {$2\mathbf q_2$};
\end{axis}
\end{tikzpicture}}{The ellipse $5x^2+4xy+8y^2=36$ and its principal axes (dashed), the eigenvector lines of the matrix with the rows $(5,2)$ and $(2,8)$. The eigenvalues $4$ and $9$ give the semi-axes $\sqrt{36/4}=3$ along $\mathbf q_1=(2,-1)/\sqrt5$ and $\sqrt{36/9}=2$ along $\mathbf q_2=(1,2)/\sqrt5$; the longer axis belongs to the smaller eigenvalue.\label{fig:ellipse-axes}}
\end{figure}

\subsection{Ellipsoids and Their Volume}

\begin{proposition}[Ellipsoids]
\label{prop:ellipsoid}
Let $A\in\mathbb R^{n\times n}$ be symmetric and positive definite, with the eigenvalues $0<\lambda_1\le\cdots\le\lambda_n$ and the orthonormal eigenvectors $\mathbf q_1,\ldots,\mathbf q_n$. The ellipsoid
\[
\mathcal E_A\coloneqq\{\mathbf x\in\mathbb R^n:\langle\mathbf x,A\mathbf x\rangle\le1\}
\]
is the image of the closed unit ball under the linear map $\mathbf u\mapsto Q\Lambda^{-1/2}\mathbf u$, where $\Lambda^{-1/2}\coloneqq\operatorname{diag}(\lambda_1^{-1/2},\ldots,\lambda_n^{-1/2})$. Its semi-axes have the lengths $1/\sqrt{\lambda_j}$, along $\mathbf q_j$, and
\[
\operatorname{vol}(\mathcal E_A)=\frac{V_n}{\sqrt{\det A}} ,
\]
where $V_n$ is the volume of the unit ball of \vcref{1}{thm:ball-volume}.
\end{proposition}

\begin{proof}
For $\mathbf x=Q\Lambda^{-1/2}\mathbf u$, \cref{thm:principal-axes} gives $\langle\mathbf x,A\mathbf x\rangle=\sum_j\lambda_j(\lambda_j^{-1/2}u_j)^2=\|\mathbf u\|^2$, and the map is invertible; so $\mathbf x\in\mathcal E_A$ if and only if $\|\mathbf u\|\le1$. The map sends the endpoints $\pm\mathbf e_j$ of the axes of the ball to $\pm\lambda_j^{-1/2}\mathbf q_j$. By \vcref{1}{thm:linear-volume}, the volume of the image of the ball is $|\det(Q\Lambda^{-1/2})|\,V_n=V_n/\sqrt{\lambda_1\cdots\lambda_n}$, since $|\det Q|=1$ (\cref{prop:det-unitary}), and $\lambda_1\cdots\lambda_n=\det A$ by \cref{thm:trace-det-eigen}.
\end{proof}

For the ellipse $x^2/a^2+y^2/b^2\le1$, the matrix is $\operatorname{diag}(1/a^2,1/b^2)$, and the proposition gives the area $\pi/\sqrt{1/(a^2b^2)}=\pi ab$ found in \vcref{1}{sec:change-variables}. The proposition explains why the determinant of a positive definite matrix measures a volume, and why a nearly singular positive definite matrix belongs to a very long ellipsoid: if $\lambda_1$ is small, the semi-axis $1/\sqrt{\lambda_1}$ is long. The ratio $\sqrt{\lambda_n/\lambda_1}$ of the longest to the shortest semi-axis is the square root of the condition number $\kappa_2(A)$ of \cref{sec:rayleigh}.

\subsection{Quadratic Functions}

A quadratic function with linear and constant terms is a quadratic form centred at a shifted point, when its matrix is invertible.

\begin{proposition}[Completing the square]
\label{prop:quadratic-centre}
Let $A\in\mathbb R^{n\times n}$ be symmetric and invertible, $\mathbf b\in\mathbb R^n$, $c\in\mathbb R$, and $f(\mathbf x)\coloneqq\langle\mathbf x,A\mathbf x\rangle+\langle\mathbf b,\mathbf x\rangle+c$. With $\mathbf x_0\coloneqq-A^{-1}\mathbf b/2$,
\[
f(\mathbf x)=\langle\mathbf x-\mathbf x_0,A(\mathbf x-\mathbf x_0)\rangle+f(\mathbf x_0),\qquad f(\mathbf x_0)=c-\tfrac14\langle\mathbf b,A^{-1}\mathbf b\rangle .
\]
The point $\mathbf x_0$ is the only critical point of $f$. If $A$ is positive definite, $f$ has its unique global minimum at $\mathbf x_0$; if $A$ is indefinite, $\mathbf x_0$ is a saddle point.
\end{proposition}

\begin{proof}
Expanding, $\langle\mathbf x-\mathbf x_0,A(\mathbf x-\mathbf x_0)\rangle=\langle\mathbf x,A\mathbf x\rangle-2\langle A\mathbf x_0,\mathbf x\rangle+\langle\mathbf x_0,A\mathbf x_0\rangle$, by the symmetry of $A$, and $-2A\mathbf x_0=\mathbf b$, $\langle\mathbf x_0,A\mathbf x_0\rangle=\langle\mathbf b,A^{-1}\mathbf b\rangle/4$; this gives the identity and the value at $\mathbf x_0$. The gradient is $\nabla f(\mathbf x)=2A\mathbf x+\mathbf b$ (\vcref{1}{sec:differentiability}), which vanishes only at $\mathbf x_0$. If $A$ is positive definite, the quadratic form in the identity is positive for $\mathbf x\ne\mathbf x_0$; if $A$ is indefinite, it takes both signs in every ball around $\mathbf x_0$, along eigenvectors for a positive and for a negative eigenvalue.
\end{proof}

The level sets of $f$ are therefore the quadrics of the previous subsection, translated to the centre $\mathbf x_0$, and minimising a quadratic function with a positive definite matrix is the same as solving the linear system $A\mathbf x=-\mathbf b/2$; \cref{sec:gaussian-uses} returns to this equivalence, with the more convenient normalisation $\langle\mathbf x,A\mathbf x\rangle/2-\langle\mathbf b,\mathbf x\rangle$. If $A$ is singular, $f$ may have no critical point at all, as $f(x,y)=x^2+y$ shows, or a whole line of them, as $f(x,y)=x^2$. The multivariate normal density of \vcref{1}{sec:gaussian} is $e^{-f}$ up to a constant factor, for $A=\Sigma^{-1}/2$ with a covariance matrix $\Sigma$; its level sets are ellipsoids centred at the mean, whose axes are the eigenvector lines of $\Sigma$ and whose semi-axes are proportional to the square roots of the eigenvalues of $\Sigma$, since $\Sigma^{-1}$ has the same eigenvectors and the reciprocal eigenvalues.

\section{Positive Definite Matrices}
\label{sec:positive-definite}

\subsection{Definitions}

\Vcref{1}{def:definite} defined definiteness for real symmetric matrices. Over $\mathbb C$ the same definition is made for Hermitian matrices, whose forms are real by \cref{prop:form-matrix}\,(iii).

\begin{definition}[Definiteness of Hermitian matrices]
\label{def:hermitian-definite}
A Hermitian matrix $A\in\mathbb F^{n\times n}$ is \term{positive definite}, written $A\succ0$, if $\mathbf x^{\mathsf H}A\mathbf x>0$ for all $\mathbf x\in\mathbb F^n\setminus\{\mathbf 0\}$, and \term{positive semidefinite}, $A\succeq0$, if $\mathbf x^{\mathsf H}A\mathbf x\ge0$ for all $\mathbf x$. It is \term{negative definite} or \term{negative semidefinite} if $-A$ is positive definite or semidefinite, and \term{indefinite} if $\mathbf x^{\mathsf H}A\mathbf x$ takes positive and negative values. For Hermitian $A$ and $B$ we write $A\succ B$ and $A\succeq B$ if $A-B\succ0$ and $A-B\succeq0$.
\end{definition}

For a real symmetric matrix it does not matter whether real or complex vectors are used. If $\mathbf x=\mathbf u+i\mathbf v$ with real $\mathbf u$ and $\mathbf v$, then
\[
\mathbf x^{\mathsf H}A\mathbf x=(\mathbf u-i\mathbf v)^{\mathsf T}A(\mathbf u+i\mathbf v)=\mathbf u^{\mathsf T}A\mathbf u+\mathbf v^{\mathsf T}A\mathbf v+i\bigl(\mathbf u^{\mathsf T}A\mathbf v-\mathbf v^{\mathsf T}A\mathbf u\bigr)=\mathbf u^{\mathsf T}A\mathbf u+\mathbf v^{\mathsf T}A\mathbf v ,
\]
since $\mathbf v^{\mathsf T}A\mathbf u=\mathbf u^{\mathsf T}A^{\mathsf T}\mathbf v=\mathbf u^{\mathsf T}A\mathbf v$; so the complex form is positive for all $\mathbf x\ne\mathbf 0$ exactly when the real form is positive for all real $\mathbf u\ne\mathbf 0$. Symmetry is part of the definition. A nonsymmetric real matrix such as the one with the rows $(1,1)$ and $(-1,1)$ has $\langle\mathbf x,A\mathbf x\rangle=\|\mathbf x\|^2>0$ for real $\mathbf x\ne\mathbf 0$, but its eigenvalues $1\pm i$ are not real, and none of the criteria below applies to it; by \cref{prop:symmetric-parts}, only its symmetric part $I$ is seen by the form.

\subsection{Criteria for Positive Definiteness}

The leading principal submatrix $A_k$ of $A$ consists of the first $k$ rows and columns, and $\Delta_k\coloneqq\det A_k$ is the $k$th leading principal minor (\cref{sec:lu,prop:pivots-minors}). The next theorem collects seven characterisations of positive definiteness. Gaussian elimination and the factorisations of \cref{sec:lu} were developed there for real matrices, but they use only the arithmetic operations, and the proofs of \cref{thm:lu-existence,thm:lu-unique,prop:ldu,prop:pivots-minors} hold word for word over $\mathbb C$; we use them in this form.

\begin{theorem}[Criteria for positive definiteness]
\label{thm:pd-criteria}
For a Hermitian matrix $A\in\mathbb F^{n\times n}$ the following are equivalent:
\begin{enumerate}
\item $A$ is positive definite;
\item all eigenvalues of $A$ are positive;
\item $A=B^{\mathsf H}B$ for some $B\in\mathbb F^{m\times n}$ with linearly independent columns;
\item all leading principal minors are positive, $\Delta_k>0$ for $k=1,\ldots,n$ (\term{Sylvester's criterion});
\item Gaussian elimination without row exchanges can be carried out on $A$, and all $n$ pivots are positive;
\item $A=LDL^{\mathsf H}$ with a unit lower triangular $L$ and a diagonal $D$ with positive diagonal entries;
\item $A=R^{\mathsf H}R$ with an upper triangular $R$ with positive diagonal entries.
\end{enumerate}
The pivots in (v) and the diagonal entries of $D$ in (vi) are the numbers $d_k=\Delta_k/\Delta_{k-1}$, $\Delta_0\coloneqq1$. Moreover, $A$ is negative definite if and only if $(-1)^k\Delta_k>0$ for $k=1,\ldots,n$.
\end{theorem}

\begin{proof}
(i)$\Leftrightarrow$(ii) is part of \cref{thm:principal-axes}.

(i)$\Rightarrow$(iv). Let $E_k\in\mathbb F^{n\times k}$ consist of the first $k$ columns of $I_n$, so that $A_k=E_k^{\mathsf H}AE_k$. For $\mathbf y\in\mathbb F^k\setminus\{\mathbf 0\}$, $E_k\mathbf y$ is $\mathbf y$ followed by zeros and is not zero, so $\mathbf y^{\mathsf H}A_k\mathbf y=(E_k\mathbf y)^{\mathsf H}A(E_k\mathbf y)>0$. Thus $A_k$ is Hermitian and positive definite, its eigenvalues are positive by (i)$\Rightarrow$(ii), and $\Delta_k$, their product (\cref{thm:trace-det-eigen}), is positive.

(iv)$\Rightarrow$(v). Since $\Delta_1,\ldots,\Delta_{n-1}$ are nonzero, the leading principal submatrices $A_1,\ldots,A_{n-1}$ are invertible, and elimination without exchanges finds nonzero pivots in the steps $1,\ldots,n-1$ (\cref{thm:lu-existence}). By \cref{prop:pivots-minors} the pivots are $u_{kk}=\Delta_k/\Delta_{k-1}>0$, the last one included.

(v)$\Rightarrow$(vi). Elimination gives $A=LU$, and \cref{prop:ldu} gives $A=LDV$ with $D=\operatorname{diag}(u_{11},\ldots,u_{nn})$ and a unit upper triangular $V$, uniquely. Taking the conjugate transpose, $A=A^{\mathsf H}=V^{\mathsf H}\bar DL^{\mathsf H}$ is another factorisation of the same kind, and uniqueness gives $V^{\mathsf H}=L$ and $\bar D=D$. So $A=LDL^{\mathsf H}$, and the pivots, which are positive, are the diagonal entries of $D$.

(vi)$\Rightarrow$(vii). Let $D^{1/2}\coloneqq\operatorname{diag}(\sqrt{d_1},\ldots,\sqrt{d_n})$ and $R\coloneqq D^{1/2}L^{\mathsf H}$. Then $R$ is upper triangular with the positive diagonal entries $\sqrt{d_k}$, and $R^{\mathsf H}R=LD^{1/2}D^{1/2}L^{\mathsf H}=LDL^{\mathsf H}=A$.

(vii)$\Rightarrow$(iii). Take $B=R$, which is invertible by \cref{thm:triangular}, so its columns are independent.

(iii)$\Rightarrow$(i). For $\mathbf x\ne\mathbf 0$, $B\mathbf x\ne\mathbf 0$ because the columns of $B$ are independent, and $\mathbf x^{\mathsf H}A\mathbf x=\langle B\mathbf x,B\mathbf x\rangle=\|B\mathbf x\|^2>0$.

The identification of the pivots was made in the proof of (iv)$\Rightarrow$(v). Finally, $A$ is negative definite if and only if $-A$ is positive definite, and the leading principal minors of $-A$ are $(-1)^k\Delta_k$.
\end{proof}

For a real symmetric matrix, (ii) and (iv) are the two parts of \vcref{1}{thm:definiteness-criteria}, which is now proved completely, and (v) is the converse announced after \cref{prop:ldu}: a symmetric matrix is positive definite exactly when elimination without exchanges produces only positive pivots. In the real case all the factors $B$, $L$, $D$ and $R$ are real.

The criteria differ greatly in cost. The eigenvalues in (ii) are the hardest to compute. The minors in (iv) are $n$ determinants, but they need not be computed separately: by \cref{prop:pivots-minors} they are the products of the pivots, so (iv) and (v) amount to the same computation, a single elimination of about $n^3/3$ operations when the symmetry is used, as \cref{sec:cholesky} shows. In practice positive definiteness is tested by attempting the factorisation (vii). Criterion (iii) is the one that most often proves positive definiteness in applications, without any computation: a Gram matrix of independent vectors is positive definite.

For the second difference matrix, $\Delta_k=k+1$, the product of the pivots $(j+1)/j$, $j\le k$, found in \cref{sec:lu}, so $T_n$ is positive definite, as its eigenvalues $2-2\cos(j\pi/(n+1))>0$ confirm. The Hilbert matrix $H_n$ of \cref{sec:pivoting}, with the entries $1/(i+j-1)=\int_0^1x^{i-1}x^{j-1}\,dx$, is the Gram matrix of the independent functions $1,x,\ldots,x^{n-1}$ in the inner product $\int_0^1fg$, hence positive definite: by \cref{lem:gram-matrix}, $\mathbf c^{\mathsf T}H_n\mathbf c=\int_0^1\bigl(\sum_jc_jx^{j-1}\bigr)^2dx$, which is positive unless the polynomial vanishes, that is, unless $\mathbf c=\mathbf 0$. (In this chapter $H_n$ denotes the Hilbert matrix, not the harmonic number of the Notation.) The matrix with the rows $(1,2)$ and $(2,1)$ has $\Delta_1=1$ and $\Delta_2=-3$ and is not positive definite; since $\Delta_2<0$ is the product of the two eigenvalues, they have opposite signs, and the matrix is indefinite. The matrix of \vcref{1}{sec:extrema-n} with the rows $(2,1,0)$, $(1,2,1)$ and $(0,1,2)$ has $\Delta_1=2$, $\Delta_2=3$ and $\Delta_3=2\cdot3-1\cdot2=4$, and is positive definite.

\subsection{Semidefinite Matrices}

Positive semidefinite matrices have the corresponding characterisations, with one important difference: the leading principal minors no longer suffice.

\begin{theorem}[Criteria for positive semidefiniteness]
\label{thm:psd-criteria}
For a Hermitian matrix $A\in\mathbb F^{n\times n}$ the following are equivalent:
\begin{enumerate}
\item $A$ is positive semidefinite;
\item all eigenvalues of $A$ are nonnegative;
\item $A=B^{\mathsf H}B$ for some $B\in\mathbb F^{m\times n}$, and $B$ can be chosen with $m=\operatorname{rank}A$;
\item all principal minors $\det A_S$, for every nonempty $S\subseteq\{1,\ldots,n\}$, are nonnegative.
\end{enumerate}
The rank of a positive semidefinite matrix is the number of its positive eigenvalues.
\end{theorem}

\begin{proof}
(i)$\Leftrightarrow$(ii) is part of \cref{thm:principal-axes}. (ii)$\Rightarrow$(iii). Let $A=Q\Lambda Q^{\mathsf H}$ with $\lambda_1\le\cdots\le\lambda_n$, and let $r$ be the number of positive eigenvalues, which are $\lambda_{n-r+1},\ldots,\lambda_n$. Let $B$ be the $r\times n$ matrix with the rows $\sqrt{\lambda_j}\,\mathbf q_j^{\mathsf H}$ for $j>n-r$. Then $B^{\mathsf H}B=\sum_{j>n-r}\lambda_j\mathbf q_j\mathbf q_j^{\mathsf H}=A$ by \cref{prop:spectral-decomposition}, since the other eigenvalues vanish. The rank of $A=Q\Lambda Q^{\mathsf H}$ is that of $\Lambda$, which is $r$ (\cref{prop:rank-product}). (iii)$\Rightarrow$(i). $\mathbf x^{\mathsf H}B^{\mathsf H}B\mathbf x=\|B\mathbf x\|^2\ge0$. (i)$\Rightarrow$(iv). As in the proof of \cref{thm:pd-criteria}, with the columns $\mathbf e_i$, $i\in S$, in place of the first $k$ columns of $I$, every principal submatrix $A_S$ is positive semidefinite; its eigenvalues are nonnegative, and so is their product $\det A_S$.

(iv)$\Rightarrow$(ii). Let $e_k$ be the sum of the principal minors of order $k$, and $e_0\coloneqq1$. By \cref{thm:charpoly-coefficients}, $\det(\lambda I-A)=\sum_{k=0}^n(-1)^ke_k\lambda^{n-k}$, and for $t>0$
\[
\det(tI+A)=(-1)^n\det(-tI-A)=(-1)^n\sum_{k=0}^n(-1)^ke_k(-t)^{n-k}=\sum_{k=0}^ne_kt^{n-k}\ge t^n>0 ,
\]
since all $e_k\ge0$. So $A+tI$ is invertible, and $-t$ is not an eigenvalue of $A$, for every $t>0$. The eigenvalues are real, so they are nonnegative.
\end{proof}

The leading principal minors are not enough for semidefiniteness. The matrix $\operatorname{diag}(0,-1)$ has $\Delta_1=0$ and $\Delta_2=0$, both nonnegative, but its form takes the value $-1$ at $\mathbf e_2$; the principal minor for $S=\{2\}$ is $-1$ and reveals it. The reason is that the passage from minors to pivots, $d_k=\Delta_k/\Delta_{k-1}$, breaks down as soon as a minor vanishes. Testing all principal minors would require $2^n-1$ determinants; in practice semidefiniteness is decided from the eigenvalues, or by elimination with symmetric pivoting, which Volume~IX describes. \Cref{tab:definiteness-criteria} compares the two theorems.

\begin{table}
\sdfloat{%
\begin{sdtable}{colspec={lll}, row{1}={font=\small\bfseries}}
\toprule
 & positive definite & positive semidefinite \\
\midrule
form & $\mathbf x^{\mathsf H}A\mathbf x>0$ for $\mathbf x\ne\mathbf 0$ & $\mathbf x^{\mathsf H}A\mathbf x\ge0$ for all $\mathbf x$ \\
eigenvalues & all $\lambda_j>0$ & all $\lambda_j\ge0$ \\
Gram form & $B^{\mathsf H}B$, independent columns & $B^{\mathsf H}B$, any $B$ \\
minors & leading ones $\Delta_k>0$ & all principal ones $\ge0$ \\
elimination & $A=LDL^{\mathsf H}$, pivots $d_k>0$ & $d_k\ge0$ whenever $A=LDL^{\mathsf H}$ \\
triangular factor & $A=R^{\mathsf H}R$, $r_{kk}>0$, unique & $A=R^{\mathsf H}R$, $r_{kk}\ge0$, not unique \\
\bottomrule
\end{sdtable}}{Criteria for positive definite and positive semidefinite Hermitian matrices (\cref{thm:pd-criteria,thm:psd-criteria} and the text after the table). A Hermitian matrix whose leading principal minors are nonnegative need not be semidefinite, as $\operatorname{diag}(0,-1)$ shows.\label{tab:definiteness-criteria}}
\end{table}

The last two rows for semidefinite matrices follow from what we have proved. If a positive semidefinite $A$ has a factorisation $A=LDL^{\mathsf H}$ with a unit lower triangular $L$, then $D=L^{-1}AL^{-\mathsf H}$ is positive semidefinite by \cref{prop:pd-properties}\,(iii) below, so its diagonal entries are nonnegative. A triangular factor always exists: write $A=B^{\mathsf H}B$ with a square $B$, for instance $B=A^{1/2}$ of \cref{thm:square-root} below, and let $B=QR$ be a QR factorisation, which Householder's method produces for every square matrix (\cref{thm:householder-qr}); after multiplying rows of $R$ and columns of $Q$ by numbers of modulus $1$, the diagonal of $R$ is nonnegative, and $A=R^{\mathsf H}Q^{\mathsf H}QR=R^{\mathsf H}R$. It is not unique: $\operatorname{diag}(0,1)=R^{\mathsf T}R$ for every $R$ with the rows $(0,c)$ and $(0,d)$ and $c^2+d^2=1$. The vectors on which a semidefinite form vanishes are exactly the vectors in its null space.

\begin{lemma}[The null space of a semidefinite matrix]
\label{lem:psd-kernel}
Let $A\in\mathbb F^{n\times n}$ be positive semidefinite. Then $\mathbf x^{\mathsf H}A\mathbf x=0$ if and only if $A\mathbf x=\mathbf 0$. Consequently a positive semidefinite matrix is positive definite if and only if it is invertible.
\end{lemma}

\begin{proof}
If $A\mathbf x=\mathbf 0$, the form vanishes. Conversely, write $A=B^{\mathsf H}B$ by \cref{thm:psd-criteria}\,(iii). If $\mathbf x^{\mathsf H}A\mathbf x=\|B\mathbf x\|^2=0$, then $B\mathbf x=\mathbf 0$ and $A\mathbf x=B^{\mathsf H}B\mathbf x=\mathbf 0$. So the form vanishes only at $\mathbf 0$ if and only if $\mathcal N(A)=\{\mathbf 0\}$, which for a square matrix means invertibility (\cref{thm:invertible}).
\end{proof}

The lemma is the fact used in \vcref{1}{sec:extrema-n}: a symmetric matrix with nonzero determinant is invertible, so if it is semidefinite, it is definite, and the second-derivative test decides at every nondegenerate critical point. For an indefinite matrix the lemma fails: the matrix with the rows $(1,0)$ and $(0,-1)$ has the form $x_1^2-x_2^2$, which vanishes at $(1,1)$, a vector that the matrix maps to $(1,-1)$.

\subsection{Properties}

\begin{proposition}[Properties of positive definite matrices]
\label{prop:pd-properties}
Let $A,B\in\mathbb F^{n\times n}$ be Hermitian and positive definite.
\begin{enumerate}
\item Every principal submatrix $A_S$ is positive definite. In particular $a_{ii}>0$, $|a_{ij}|^2<a_{ii}a_{jj}$ for $i\ne j$, and an entry of largest modulus lies on the diagonal.
\item $A+B$ and $cA$, $c>0$, are positive definite, and so is $A+C$ for every positive semidefinite $C$.
\item $S^{\mathsf H}AS$ is positive definite for every $S\in\mathbb F^{n\times m}$ with independent columns, and positive semidefinite for every $S\in\mathbb F^{n\times m}$; the second statement holds for every positive semidefinite $A$. In particular congruence preserves positive definiteness and positive semidefiniteness.
\item $A^{-1}$ is positive definite.
\item $\operatorname{tr}A>0$ and $0<\det A\le a_{11}a_{22}\cdots a_{nn}$.
\end{enumerate}
\end{proposition}

\begin{proof}
(i) was shown in the proof of \cref{thm:pd-criteria} for leading submatrices, and the same argument applies to every $S$. For $S=\{i\}$ it gives $a_{ii}>0$, and for $S=\{i,j\}$ the determinant $a_{ii}a_{jj}-|a_{ij}|^2$ is positive. Hence $|a_{ij}|<\sqrt{a_{ii}a_{jj}}\le\max(a_{ii},a_{jj})$. (ii) The forms add, and $\mathbf x^{\mathsf H}(A+C)\mathbf x\ge\mathbf x^{\mathsf H}A\mathbf x$. (iii) $\mathbf x^{\mathsf H}S^{\mathsf H}AS\mathbf x=(S\mathbf x)^{\mathsf H}A(S\mathbf x)$, which is positive if $S\mathbf x\ne\mathbf 0$, and nonnegative in any case; the second conclusion uses only $A\succeq0$. If $S$ is square and invertible, the same computation with $S^{-1}$ in place of $S$ gives the converse. (iv) $A^{-1}=A^{-1}AA^{-1}=(A^{-1})^{\mathsf H}A(A^{-1})$, since $(A^{-1})^{\mathsf H}=(A^{\mathsf H})^{-1}=A^{-1}$, so (iii) applies with $S=A^{-1}$. Equivalently, the eigenvalues of $A^{-1}$ are the numbers $1/\lambda_j>0$. (v) The trace and the determinant are the sum and the product of the eigenvalues (\cref{thm:trace-det-eigen}). For the upper bound write $A=R^{\mathsf H}R$ by \cref{thm:pd-criteria}\,(vii). The $j$th diagonal entry of $R^{\mathsf H}R$ is $\|\mathbf r_j\|^2$, the squared length of the $j$th column of $R$, and Hadamard's inequality (\cref{thm:hadamard-complex}) gives $\det A=|\det R|^2\le\prod_j\|\mathbf r_j\|^2=\prod_ja_{jj}$.
\end{proof}

Part (i) gives quick tests for the failure of positive definiteness: a symmetric matrix with a nonpositive diagonal entry, or with an off-diagonal entry larger in modulus than every diagonal entry, is not positive definite. The matrix with the rows $(1,2)$ and $(2,5)$ passes the test, $|a_{12}|^2=4<5$, and is positive definite, with $\Delta_2=1$. The bound (v) is Hadamard's inequality for positive definite matrices; it is an equality exactly for diagonal matrices, by the equality case of \cref{thm:hadamard-complex}.

\subsection{Schur Complements and Gram Matrices}

\Cref{thm:schur-complement} factorised a block matrix through the Schur complement $S=D-CA^{-1}B$. For Hermitian matrices the factorisation is a congruence, and it turns positive definiteness of a large matrix into that of two smaller ones.

\begin{theorem}[Schur complements of positive definite matrices]
\label{thm:schur-pd}
Let
\[
M=\begin{bmatrix}A&B\\B^{\mathsf H}&D\end{bmatrix}
\]
be Hermitian, with $A\in\mathbb F^{k\times k}$ and $D\in\mathbb F^{l\times l}$.
\begin{enumerate}
\item $M$ is positive definite if and only if $A$ is positive definite and the Schur complement $S\coloneqq D-B^{\mathsf H}A^{-1}B$ is positive definite.
\item If $A$ is positive definite, $M$ is positive semidefinite if and only if $S$ is positive semidefinite.
\end{enumerate}
\end{theorem}

\begin{proof}
If $M$ is positive definite, so is its leading principal submatrix $A$, by \cref{prop:pd-properties}\,(i). So in either part we may assume that $A$ is positive definite, hence invertible and Hermitian. With $C=B^{\mathsf H}$, the factorisation of \cref{thm:schur-complement}\,(i) reads
\[
M=L\begin{bmatrix}A&0\\0&S\end{bmatrix}L^{\mathsf H},\qquad L\coloneqq\begin{bmatrix}I&0\\B^{\mathsf H}A^{-1}&I\end{bmatrix},
\]
because the third factor there, with the blocks $I$, $A^{-1}B$, $0$ and $I$, is $L^{\mathsf H}$: $(B^{\mathsf H}A^{-1})^{\mathsf H}=(A^{-1})^{\mathsf H}B=A^{-1}B$. The matrix $L$ is invertible, so $M$ is congruent to $\operatorname{diag}(A,S)$, and by \cref{prop:pd-properties}\,(iii) the one is positive definite or semidefinite if and only if the other is. For $\mathbf z=(\mathbf x,\mathbf y)$, $\mathbf z^{\mathsf H}\operatorname{diag}(A,S)\mathbf z=\mathbf x^{\mathsf H}A\mathbf x+\mathbf y^{\mathsf H}S\mathbf y$, which is positive for all $\mathbf z\ne\mathbf 0$ if and only if $A$ and $S$ are positive definite, and nonnegative for all $\mathbf z$, when $A$ is positive definite, if and only if $S$ is positive semidefinite.
\end{proof}

This is the statement promised after \cref{thm:schur-complement}. Since one step of elimination leaves the Schur complement of the pivot in the remaining rows and columns, the theorem with $k=1$ is the inductive step of criterion (v): $M$ is positive definite if and only if its first pivot $a_{11}$ is positive and the matrix that elimination leaves behind is positive definite. For the matrix with the rows $(2,1)$ and $(1,2)$, the Schur complement of $a_{11}=2$ is $2-1/2=3/2$, which is the second pivot.

The Gram matrices of \cref{prop:gram,lem:gram-matrix} are the matrices of criterion (iii). For vectors $\mathbf b_1,\ldots,\mathbf b_k$ of an inner product space, the Gram matrix $G=(\langle\mathbf b_i,\mathbf b_j\rangle)$ satisfies $\mathbf c^{\mathsf H}G\mathbf c=\|\sum_jc_j\mathbf b_j\|^2$, so it is positive semidefinite, and positive definite exactly when the vectors are independent; conversely, by \cref{thm:psd-criteria}\,(iii), every positive semidefinite matrix is a Gram matrix, of the rows of $B$ taken as vectors of $\mathbb F^m$. Positive semidefinite matrices are therefore exactly the tables of inner products of families of vectors, and positive definite ones those of independent families.

\subsection{Inner Products on \texorpdfstring{$\mathbb F^n$}{Fn}}

\Cref{ch:inner-products} observed in \eqref{eq:gram-form} that every inner product on $\mathbb F^n$ has the form $\mathbf x^{\mathsf H}G\mathbf y$ with $G$ Hermitian and $\mathbf x^{\mathsf H}G\mathbf x>0$ for $\mathbf x\ne\mathbf 0$. In the language of this chapter:

\begin{corollary}[Inner products on $\mathbb F^n$]
\label{cor:inner-products-pd}
The formula $\langle\mathbf x,\mathbf y\rangle_G\coloneqq\mathbf x^{\mathsf H}G\mathbf y$ defines an inner product on $\mathbb F^n$ if and only if $G$ is Hermitian and positive definite, and every inner product on $\mathbb F^n$ is of this form. If $G=R^{\mathsf H}R$ is the factorisation of \cref{thm:pd-criteria}\,(vii), then
\[
\langle\mathbf x,\mathbf y\rangle_G=\langle R\mathbf x,R\mathbf y\rangle :
\]
every inner product on $\mathbb F^n$ is the standard inner product in the coordinates $R\mathbf x$, and $\mathbf x\mapsto R\mathbf x$ maps a basis that is orthonormal for $\langle\cdot,\cdot\rangle_G$ to one that is orthonormal for the standard inner product.
\end{corollary}

\begin{proof}
The first two statements are \eqref{eq:gram-form} and the remark after it, together with \cref{def:hermitian-definite}. Then $\mathbf x^{\mathsf H}R^{\mathsf H}R\mathbf y=(R\mathbf x)^{\mathsf H}(R\mathbf y)$, and $\langle R\mathbf u,R\mathbf v\rangle=\langle\mathbf u,\mathbf v\rangle_G$ for all vectors of a basis.
\end{proof}

The weighted inner products of \cref{sec:inner-product-spaces} are the case $G=W=\operatorname{diag}(w_i)$, with $R=W^{1/2}$. A non-diagonal example is $G$ with the rows $(2,1)$ and $(1,2)$, so that $\langle\mathbf x,\mathbf y\rangle_G=2\bar x_1y_1+\bar x_1y_2+\bar x_2y_1+2\bar x_2y_2$: it is an inner product because $\Delta_1=2$ and $\Delta_2=3$ are positive, which is the test that \cref{ch:inner-products} called less obvious than for weights.

\subsection{The Loewner Order}

The relation $A\succeq B$ of \cref{def:hermitian-definite} is the \term{Loewner order} on Hermitian matrices~\cite{loewner1934}. It is a partial order: $A\succeq A$; if $A\succeq B$ and $B\succeq C$, then $A-C=(A-B)+(B-C)\succeq0$; and if $A\succeq B$ and $B\succeq A$, then the form of $A-B$ vanishes identically, so $A=B$ by \cref{prop:form-matrix}\,(ii). It is not a total order: neither of $\operatorname{diag}(1,0)$ and $\operatorname{diag}(0,1)$ dominates the other, since their difference $\operatorname{diag}(1,-1)$ is indefinite. Congruence preserves it, by \cref{prop:pd-properties}\,(iii): $A\succeq B$ implies $S^{\mathsf H}AS\succeq S^{\mathsf H}BS$. By \cref{thm:principal-axes},
\[
\lambda_{\min}(A)\,I\preceq A\preceq\lambda_{\max}(A)\,I ,
\]
and the eigenvalues are the best constants in these inequalities. \Cref{sec:min-max} shows that $A\succeq B$ implies $\lambda_k(A)\ge\lambda_k(B)$ for every $k$; the converse is false, as the two matrices above show, which have the same eigenvalues.

\begin{note}{The cone of positive semidefinite matrices}
The positive semidefinite $n\times n$ matrices form a \emph{convex cone}: sums and nonnegative multiples of positive semidefinite matrices are positive semidefinite, by \cref{prop:pd-properties}\,(ii). The cone is closed, because each condition $\mathbf x^{\mathsf H}A\mathbf x\ge0$ is preserved by limits, and its interior consists of the positive definite matrices. Optimisation over this cone, minimising a linear function of a symmetric matrix $X$ subject to linear equations and the constraint $X\succeq0$, is called \emph{semidefinite programming}; it contains linear programming as the case of diagonal matrices, and it has become a standard tool for approximating hard combinatorial problems and for control theory~\cite{vandenberghe1996}. A further property of the cone is the theorem of Schur~\cite{schur1911}: if $A$ and $B$ are positive semidefinite, so is their entrywise product $A\circ B=(a_{ij}b_{ij})$. Indeed, write $B=\sum_k\mu_k\mathbf v_k\mathbf v_k^{\mathsf H}$ with $\mu_k\ge0$ by \cref{prop:spectral-decomposition}; the entrywise product of $A$ with $\mathbf v\mathbf v^{\mathsf H}$ has the entries $a_{ij}v_i\bar v_j$, so it is $D_{\mathbf v}AD_{\mathbf v}^{\mathsf H}$ with $D_{\mathbf v}=\operatorname{diag}(v_1,\ldots,v_n)$, which is positive semidefinite by \cref{prop:pd-properties}\,(iii), and $A\circ B$ is a nonnegative combination of such matrices~\cite[Sect.~7.5]{horn2013}. Volume~XIII treats semidefinite programming as part of convex optimisation.
\end{note}

\section{The Cholesky Factorisation}
\label{sec:cholesky}

\subsection{Existence and Uniqueness}

Criterion (vii) of \cref{thm:pd-criteria} is the factorisation that is used in computation.

\begin{theorem}[Cholesky factorisation]
\label{thm:cholesky}
Every Hermitian positive definite matrix $A\in\mathbb F^{n\times n}$ can be written in exactly one way as
\[
A=R^{\mathsf H}R
\]
with an upper triangular $R\in\mathbb F^{n\times n}$ whose diagonal entries are positive. If $A$ is real, so is $R$. The factor is $R=D^{1/2}L^{\mathsf H}$ in terms of the factorisation $A=LDL^{\mathsf H}$ of \cref{thm:pd-criteria}\,(vi), and $r_{kk}^2=\Delta_k/\Delta_{k-1}$.
\end{theorem}

\begin{proof}
Existence and the formula for $R$ are (vi)$\Rightarrow$(vii) in the proof of \cref{thm:pd-criteria}, and $r_{kk}^2=d_k$; the elimination works in real arithmetic for a real $A$. For uniqueness, let $A=R_1^{\mathsf H}R_1=R_2^{\mathsf H}R_2$. The matrix $T\coloneqq R_2R_1^{-1}$ is upper triangular with positive diagonal entries (\cref{thm:triangular}), and
\[
T^{\mathsf H}T=R_1^{-\mathsf H}R_2^{\mathsf H}R_2R_1^{-1}=R_1^{-\mathsf H}R_1^{\mathsf H}R_1R_1^{-1}=I ,
\]
where $R_1^{-\mathsf H}$ denotes $(R_1^{-1})^{\mathsf H}$. As in the uniqueness part of \cref{thm:qr}, an upper triangular matrix with positive diagonal entries and $T^{\mathsf H}T=I$ is the identity, so $R_2=R_1$.
\end{proof}

The factorisation is often written $A=GG^{\mathsf H}$ with the lower triangular factor $G=R^{\mathsf H}$, as in \cref{sec:special-systems}; the two forms are the same factorisation. The lower triangular form is the one returned by the routine \code{linalg.cholesky} of the stardust library, unless the upper one is requested. For a $2\times2$ matrix with the rows $(a,b)$ and $(\bar b,c)$,
\[
\begin{bmatrix}a&b\\\bar b&c\end{bmatrix}=\begin{bmatrix}\sqrt a&0\\\bar b/\sqrt a&\sqrt{c-|b|^2/a}\end{bmatrix}\begin{bmatrix}\sqrt a&b/\sqrt a\\0&\sqrt{c-|b|^2/a}\end{bmatrix},
\]
which is completing the square once more, and for the matrix with the rows $(4,2,-2)$, $(2,5,1)$ and $(-2,1,6)$,
\[
\begin{bmatrix}4&2&-2\\2&5&1\\-2&1&6\end{bmatrix}=\begin{bmatrix}2&0&0\\1&2&0\\-1&1&2\end{bmatrix}\begin{bmatrix}2&1&-1\\0&2&1\\0&0&2\end{bmatrix},
\]
as one checks by multiplying out; the algorithm below produces the factor $R$ on the right column by column.

\subsection{The Algorithm}

Comparing the entry $(i,j)$, $i\le j$, on both sides of $A=R^{\mathsf H}R$ gives
\[
a_{ij}=\sum_{k=1}^{i}\bar r_{ki}r_{kj},
\]
since $r_{ki}=0$ for $k>i$. For $i<j$ the last term is $\bar r_{ii}r_{ij}=r_{ii}r_{ij}$, and for $i=j$ it is $r_{jj}^2$. Solving for these terms,
\begin{equation}
r_{ij}=\frac1{r_{ii}}\Bigl(a_{ij}-\sum_{k=1}^{i-1}\bar r_{ki}r_{kj}\Bigr)\quad(i<j),\qquad r_{jj}=\Bigl(a_{jj}-\sum_{k=1}^{j-1}|r_{kj}|^2\Bigr)^{1/2},
\label{eq:cholesky-entries}
\end{equation}
and the entries can be computed column by column, from the top: every quantity on the right has been computed before it is needed. This is \cref{alg:cholesky}. The algorithm reads only the upper triangle of $A$, and it can overwrite it with $R$.

\begin{algorithm}
\begin{algorithmic}[1]
\Require a Hermitian matrix $A\in\mathbb F^{n\times n}$ (only $a_{ij}$ with $i\le j$ are read)
\Ensure the upper triangular $R$ with $A=R^{\mathsf H}R$, or the report that $A$ is not positive definite
\Name{Cholesky}{A}
\For{$j=1,\ldots,n$}
  \For{$i=1,\ldots,j-1$}
    \State $r_{ij}\gets\bigl(a_{ij}-\sum_{k=1}^{i-1}\bar r_{ki}r_{kj}\bigr)/r_{ii}$
  \EndFor
  \State $s\gets a_{jj}-\sum_{k=1}^{j-1}|r_{kj}|^2$ \Comment{$s=\Delta_j/\Delta_{j-1}$}
  \If{$s\le0$}
    \State \Return ``not positive definite'' \Comment{$A_j$ is not positive definite}
  \EndIf
  \State $r_{jj}\gets\sqrt s$
\EndFor
\State \Return $R$
\end{algorithmic}
\caption{The Cholesky factorisation, column by column (\cref{thm:cholesky,prop:cholesky-test}). It uses $n^3/3+n^2/2-5n/6$ arithmetic operations and $n$ square roots for a real matrix; the entries below the diagonal of $R$ are zero and are neither computed nor stored.}
\label{alg:cholesky}
\end{algorithm}

\begin{proposition}[The cost of the Cholesky factorisation]
\label{prop:cholesky-cost}
For a real matrix, \cref{alg:cholesky} uses $n^3/3+n^2/2-5n/6$ additions, subtractions, multiplications and divisions, and $n$ square roots.
\end{proposition}

\begin{proof}
For fixed $j$ and $i<j$, line~3 uses $i-1$ multiplications, $i-1$ subtractions and one division, $2i-1$ operations in all, and the sum over $i=1,\ldots,j-1$ is $(j-1)^2$. Line~4 uses $j-1$ multiplications and $j-1$ subtractions. With $m=j-1$, column $j$ costs $m^2+2m$ operations and one square root, and
\[
\sum_{m=0}^{n-1}(m^2+2m)=\frac{(n-1)n(2n-1)}6+n(n-1)=\frac{n^3}3+\frac{n^2}2-\frac{5n}6 .\qedhere
\]
\end{proof}

The leading term $n^3/3$ is half the $2n^3/3$ of Gaussian elimination (\cref{prop:elimination-cost}): the symmetry halves the number of entries to be computed, and the factorisation needs no pivoting. That pivoting is not needed for stability either is a theorem of Volume~IX: the Cholesky factorisation computed in floating-point arithmetic is the exact factorisation of a matrix $A+E$ with $\|E\|$ of the order of $u\|A\|$, without any exchanges of rows~\cite[Chap.~10]{higham2002accuracy}. The reason, in short, is that the entries of $R$ cannot grow, since $|r_{kj}|^2\le\sum_k|r_{kj}|^2=a_{jj}$ by \eqref{eq:cholesky-entries}. For large sparse matrices, such as those of discretised differential equations, the factor $R$ has nonzero entries in many positions where $A$ has zeros, and Volume~X treats the orderings of the unknowns that limit this fill.

\subsection{A Test for Positive Definiteness}

\Cref{alg:cholesky} stops when it meets a nonpositive number under the square root. This happens exactly when the matrix is not positive definite, so the algorithm doubles as the cheapest test.

\begin{proposition}[Cholesky as a test]
\label{prop:cholesky-test}
Let $A\in\mathbb F^{n\times n}$ be Hermitian. \Cref{alg:cholesky} runs to the end if and only if $A$ is positive definite, and then it returns the factor of \cref{thm:cholesky}. If it stops in column $j$, then $A_1,\ldots,A_{j-1}$ are positive definite and $A_j$ is not. In every column $j$ that it reaches, $s=\Delta_j/\Delta_{j-1}$.
\end{proposition}

\begin{proof}
Suppose that the columns $1,\ldots,j-1$ have been computed with positive numbers $s$, so that the upper triangular matrix $R_{j-1}$ formed by them has positive diagonal entries and, by \eqref{eq:cholesky-entries}, $A_{j-1}=R_{j-1}^{\mathsf H}R_{j-1}$. In column $j$, line~3 computes the vector $\mathbf r=(r_{1j},\ldots,r_{j-1,j})$ that solves $R_{j-1}^{\mathsf H}\mathbf r=\mathbf a$, where $\mathbf a=(a_{1j},\ldots,a_{j-1,j})$, by forward substitution, and line~4 computes $s=a_{jj}-\|\mathbf r\|^2$. By block multiplication,
\[
A_j=\begin{bmatrix}A_{j-1}&\mathbf a\\\mathbf a^{\mathsf H}&a_{jj}\end{bmatrix}=\begin{bmatrix}R_{j-1}^{\mathsf H}&\mathbf 0\\\mathbf r^{\mathsf H}&1\end{bmatrix}\begin{bmatrix}I&\mathbf 0\\\mathbf 0^{\mathsf T}&s\end{bmatrix}\begin{bmatrix}R_{j-1}&\mathbf r\\\mathbf 0^{\mathsf T}&1\end{bmatrix}.
\]
The outer factors are invertible and conjugate transposes of each other, so $A_j$ is congruent to $\operatorname{diag}(I_{j-1},s)$, and it is positive definite if and only if $s>0$ (\cref{prop:pd-properties}\,(iii)). Taking determinants, $\Delta_j=|\det R_{j-1}|^2s=\Delta_{j-1}s$. By induction on $j$, the algorithm reaches column $j$ and passes it if and only if $A_1,\ldots,A_j$ are positive definite; it runs to the end if and only if $A=A_n$ is positive definite, and then the computed $R$ satisfies $A=R^{\mathsf H}R$ and is the factor of \cref{thm:cholesky}.
\end{proof}

For the matrix with the rows $(1,2)$ and $(2,1)$, the first column gives $r_{11}=1$, and the second $r_{12}=2$ and $s=1-4=-3<0$: the algorithm stops, and the matrix is not positive definite. In floating-point arithmetic the test is reliable only up to the rounding errors: a positive definite matrix whose smallest eigenvalue is tiny compared with its norm may be reported as not positive definite, and \cref{sec:symmetric-computation} shows this with the Hilbert matrices.

\subsection{Normal Equations and the QR Factorisation}

If $A\in\mathbb F^{m\times n}$ has independent columns, $A^{\mathsf H}A$ is positive definite, and the QR factorisation of $A$ yields its Cholesky factor.

\begin{proposition}[Cholesky and QR]
\label{prop:cholesky-qr}
Let $A\in\mathbb F^{m\times n}$ have linearly independent columns and the QR factorisation $A=QR$ of \cref{thm:qr}. Then $R$ is the Cholesky factor of $A^{\mathsf H}A$: $A^{\mathsf H}A=R^{\mathsf H}R$. The least squares solution of $A\mathbf x\approx\mathbf b$ is obtained from the normal equations by solving $R^{\mathsf H}\mathbf y=A^{\mathsf H}\mathbf b$ and $R\hat{\mathbf x}=\mathbf y$.
\end{proposition}

\begin{proof}
$A^{\mathsf H}A=R^{\mathsf H}Q^{\mathsf H}QR=R^{\mathsf H}R$, and $R$ is upper triangular with positive diagonal entries, so it is the Cholesky factor by the uniqueness in \cref{thm:cholesky}. The normal equations $R^{\mathsf H}R\hat{\mathbf x}=A^{\mathsf H}\mathbf b$ (\cref{thm:least-squares}) split into the two triangular systems.
\end{proof}

This is the statement promised after \cref{thm:qr}. The two routes to $R$ differ in cost and in accuracy. Forming $A^{\mathsf H}A$ takes about $mn^2$ operations for a real matrix, when only one triangle is computed, and its Cholesky factorisation $n^3/3$ more; the Householder QR factorisation takes $2mn^2-2n^3/3$ (\cref{prop:householder-cost}). But $\kappa(A^{\mathsf H}A)=\kappa(A)^2$, as the L\"auchli example of \cref{sec:least-squares} showed, and the Cholesky factor computed from $A^{\mathsf H}A$ inherits the larger condition number: in that example, for $\varepsilon=10^{-8}$, the computed $A^{\mathsf T}A$ was singular and \code{linalg.cholesky} stopped, which by \cref{prop:cholesky-test} it had to do for a matrix that is not positive definite. The normal equations solved by Cholesky are the method of choice when $A$ is well conditioned and $m\gg n$, and the QR factorisation when it is not.

\subsection{The Square Root of a Positive Semidefinite Matrix}

A positive number has a unique positive square root, and the spectral theorem extends this to matrices.

\begin{theorem}[Square root]
\label{thm:square-root}
Let $A\in\mathbb F^{n\times n}$ be positive semidefinite, $A=Q\Lambda Q^{\mathsf H}$. Then
\[
A^{1/2}\coloneqq Q\Lambda^{1/2}Q^{\mathsf H},\qquad\Lambda^{1/2}=\operatorname{diag}(\sqrt{\lambda_1},\ldots,\sqrt{\lambda_n}),
\]
is the only positive semidefinite matrix whose square is $A$. It is positive definite if $A$ is, and it is a polynomial in $A$, so that it commutes with every matrix that commutes with $A$.
\end{theorem}

\begin{proof}
$A^{1/2}$ is Hermitian with the nonnegative eigenvalues $\sqrt{\lambda_j}$, hence positive semidefinite, positive definite if all $\lambda_j>0$, and $(A^{1/2})^2=Q\Lambda^{1/2}Q^{\mathsf H}Q\Lambda^{1/2}Q^{\mathsf H}=Q\Lambda Q^{\mathsf H}=A$. Let $\mu_1,\ldots,\mu_k$ be the different eigenvalues of $A$ and $P_i$ the orthogonal projections onto their eigenspaces, so that $A^{1/2}=\sum_i\sqrt{\mu_i}P_i$ by \cref{prop:spectral-decomposition}. With the Lagrange polynomial $p$ of degree less than $k$ that takes the values $p(\mu_i)=\sqrt{\mu_i}$ (\cref{thm:interpolation}), $p(A)=\sum_ip(\mu_i)P_i=A^{1/2}$; if $BA=AB$, then $B$ commutes with every power of $A$ and so with $p(A)$. For uniqueness, let $C\succeq0$ with $C^2=A$, and let $C=\sum_l\nu_lP'_l$ be its spectral decomposition, with pairwise different $\nu_l\ge0$. Then $A=C^2=\sum_l\nu_l^2P'_l$ by \cref{prop:spectral-decomposition}, applied to $C$ with $p(x)=x^2$, and the numbers $\nu_l^2$ are pairwise different, because the $\nu_l$ are nonnegative. Let $\mathbf v\in\mathbb F^n$ and $\mathbf v_l\coloneqq P'_l\mathbf v$, so that $\mathbf v=\sum_l\mathbf v_l$, since $\sum_lP'_l=I$, and $A\mathbf v=\sum_l\nu_l^2\mathbf v_l$, since $P'_mP'_l=0$ for $m\ne l$ and $P'_lP'_l=P'_l$. If $A\mathbf v=\mu\mathbf v$, then $\sum_l(\nu_l^2-\mu)\mathbf v_l=\mathbf 0$, and since the vectors $\mathbf v_l$ lie in the mutually orthogonal images of the $P'_l$, each term vanishes: $\mathbf v_l=\mathbf 0$ whenever $\nu_l^2\ne\mu$. Hence every eigenvalue $\mu$ of $A$ is one of the numbers $\nu_l^2$, and its eigenspace is contained in the image of the $P'_l$ with $\nu_l^2=\mu$; conversely, every nonzero vector of that image is an eigenvector of $A$ for $\nu_l^2$. So the eigenvalues of $A$ are exactly the numbers $\nu_l^2$, the eigenspace of $\nu_l^2$ is the image of $P'_l$, and $P'_l$, being the orthogonal projection onto this image, is the projection $P_i$ for $\mu_i=\nu_l^2$. Hence $C=\sum_l\nu_lP'_l=\sum_i\sqrt{\mu_i}P_i=A^{1/2}$.
\end{proof}

The square root and the Cholesky factor both satisfy $A=X^{\mathsf H}X$, but they are different matrices. For
\[
A=\begin{bmatrix}5&4\\4&5\end{bmatrix},\qquad A^{1/2}=\begin{bmatrix}2&1\\1&2\end{bmatrix},\qquad R=\frac1{\sqrt5}\begin{bmatrix}5&4\\0&3\end{bmatrix},
\]
where $A$ has the eigenvalues $9$ and $1$ with the eigenvectors $(1,1)$ and $(1,-1)$, so that $A^{1/2}$ has the eigenvalues $3$ and $1$ with the same eigenvectors; $A^{1/2}$ is symmetric and $R$ is triangular. In general, if $X^{\mathsf H}X=A$ for two invertible matrices $X=X_1$ and $X=X_2$, then $U\coloneqq X_2X_1^{-1}$ satisfies $U^{\mathsf H}U=X_1^{-\mathsf H}AX_1^{-1}=I$, so all square factors of a positive definite $A$ are $UA^{1/2}$ with unitary $U$; in particular $R=UA^{1/2}$ for a unitary $U$. \Cref{ch:svd} meets this factorisation of $R$ as its polar decomposition. The Cholesky factor is cheap, a finite computation of about $n^3/3$ operations, and the square root needs the eigenvalues; the square root is the natural object in theory, where its symmetry and its commutation with $A$ matter, and the Cholesky factor in computation.

The square root completes a property of the Loewner order: if $A\succeq B\succ0$, then $B^{-1}\succeq A^{-1}$. Indeed, by congruence with $B^{-1/2}$, the inverse of $B^{1/2}$, the hypothesis becomes $C\coloneqq B^{-1/2}AB^{-1/2}\succeq I$, so all eigenvalues of $C$ are at least $1$; then the eigenvalues of $C^{-1}=B^{1/2}A^{-1}B^{1/2}$ are at most $1$, that is $C^{-1}\preceq I$, and congruence with $B^{-1/2}$ gives $A^{-1}\preceq B^{-1}$.

\subsection{Weighted Least Squares}

\Cref{sec:least-squares} weighted the measurements of a least squares problem by positive numbers. When the errors of the measurements are correlated, the weights become a positive definite matrix $W$ (Volume~XVI), and the quantity to be minimised is $(\mathbf b-A\mathbf x)^{\mathsf H}W(\mathbf b-A\mathbf x)$, the squared length of the residual in the inner product $\langle\cdot,\cdot\rangle_W$ of \cref{cor:inner-products-pd}.

\begin{theorem}[Weighted least squares]
\label{thm:weighted-ls}
Let $A\in\mathbb F^{m\times n}$, $\mathbf b\in\mathbb F^m$, and let $W\in\mathbb F^{m\times m}$ be Hermitian and positive definite with the Cholesky factorisation $W=R^{\mathsf H}R$. For $\hat{\mathbf x}\in\mathbb F^n$ the following are equivalent:
\begin{enumerate}
\item $\hat{\mathbf x}$ minimises $\|\mathbf b-A\mathbf x\|_W^2\coloneqq(\mathbf b-A\mathbf x)^{\mathsf H}W(\mathbf b-A\mathbf x)$ over $\mathbf x\in\mathbb F^n$;
\item $\hat{\mathbf x}$ is a least squares solution of $RA\mathbf x\approx R\mathbf b$ in the standard inner product;
\item $\hat{\mathbf x}$ solves the weighted normal equations $A^{\mathsf H}WA\hat{\mathbf x}=A^{\mathsf H}W\mathbf b$.
\end{enumerate}
Such $\hat{\mathbf x}$ exist, and they are unique if and only if the columns of $A$ are linearly independent; $A^{\mathsf H}WA$ is then positive definite.
\end{theorem}

\begin{proof}
By \cref{cor:inner-products-pd}, $\|\mathbf b-A\mathbf x\|_W^2=\|R\mathbf b-RA\mathbf x\|^2$, which gives (i)$\Leftrightarrow$(ii). By \cref{thm:least-squares}, (ii) is equivalent to the normal equations $(RA)^{\mathsf H}RA\hat{\mathbf x}=(RA)^{\mathsf H}R\mathbf b$, which are (iii) since $R^{\mathsf H}R=W$; the same theorem gives existence, and uniqueness if and only if the columns of $RA$ are independent. Since $R$ is invertible, $RA\mathbf x=\mathbf 0$ only if $A\mathbf x=\mathbf 0$, so the columns of $RA$ are independent exactly when those of $A$ are, and then $A^{\mathsf H}WA=(RA)^{\mathsf H}(RA)$ is positive definite by \cref{thm:pd-criteria}\,(iii).
\end{proof}

For a diagonal $W$ the theorem is the weighted problem of \cref{sec:least-squares}, with $R=W^{1/2}$. In computation one does not form $A^{\mathsf H}WA$ but computes the Cholesky factor $R$ of $W$, forms $RA$ and $R\mathbf b$ by triangular products, and solves the ordinary least squares problem by the QR factorisation of $RA$. The case of a general $W$ is the \term{generalised least squares} problem of statistics, where $W$ is the inverse of the covariance matrix of the errors.

\section{The Rayleigh Quotient}
\label{sec:rayleigh}

\subsection{Extreme Values}

The \term{Rayleigh quotient} of a Hermitian matrix $A\in\mathbb F^{n\times n}$ is the function
\[
r_A(\mathbf x)\coloneqq\frac{\mathbf x^{\mathsf H}A\mathbf x}{\mathbf x^{\mathsf H}\mathbf x}=\frac{q_A(\mathbf x)}{\|\mathbf x\|^2}\qquad(\mathbf x\ne\mathbf 0),
\]
which \cref{sec:computing-eigenvalues} used to estimate eigenvalues. It is real, and it does not change when $\mathbf x$ is multiplied by a nonzero number, so it is a function of the direction of $\mathbf x$ alone, and its values are the values of $q_A$ on the unit sphere. For an eigenvector, $r_A(\mathbf x)=\lambda$.

\begin{theorem}[Rayleigh quotient]
\label{thm:rayleigh}
Let $A\in\mathbb F^{n\times n}$ be Hermitian. Then
\[
\lambda_1=\min_{\mathbf x\ne\mathbf 0}r_A(\mathbf x),\qquad\lambda_n=\max_{\mathbf x\ne\mathbf 0}r_A(\mathbf x),
\]
and the set of values of $r_A$ is the interval $[\lambda_1,\lambda_n]$. The maximum is attained exactly at the eigenvectors for $\lambda_n$, and the minimum exactly at the eigenvectors for $\lambda_1$.
\end{theorem}

\begin{proof}
By \cref{thm:principal-axes}, with $\mathbf y=Q^{\mathsf H}\mathbf x$,
\[
r_A(\mathbf x)=\frac{\sum_j\lambda_j|y_j|^2}{\sum_j|y_j|^2},
\]
a weighted average of the eigenvalues with the nonnegative weights $|y_j|^2/\|\mathbf y\|^2$, which add up to $1$. So $\lambda_1\le r_A(\mathbf x)\le\lambda_n$, with $r_A(\mathbf q_1)=\lambda_1$ and $r_A(\mathbf q_n)=\lambda_n$. The equality $r_A(\mathbf x)=\lambda_n$ holds exactly when all the weight lies on indices with $\lambda_j=\lambda_n$, that is, when $\mathbf x=Q\mathbf y$ is a combination of the eigenvectors $\mathbf q_j$ with $\lambda_j=\lambda_n$, which means $\mathbf x\in E_{\lambda_n}$; similarly for $\lambda_1$. The function $t\mapsto r_A\bigl((1-t)\mathbf q_1+t\mathbf q_n\bigr)$ is continuous on $[0,1]$, its argument never vanishes since the two vectors are orthonormal, and it runs from $\lambda_1$ to $\lambda_n$; by the intermediate value theorem every value in between is attained.
\end{proof}

This is the statement found with Lagrange multipliers in \vcref{1}{sec:lagrange}, now with the information on where the extremes are attained. It has an immediate consequence for the second-derivative test: the constant $\mu$ of \vcref{1}{lem:definite-bound} can be taken to be $\lambda_1(A)$, and no larger constant works, since $\langle\mathbf h,A\mathbf h\rangle=\lambda_1\|\mathbf h\|^2$ for an eigenvector. At a critical point of a $C^2$ function the Hessian therefore gives a strict local minimum when $\lambda_{\min}(\nabla^2f(\mathbf c))>0$, a strict local maximum when $\lambda_{\max}(\nabla^2f(\mathbf c))<0$, and a saddle point when $\lambda_{\min}<0<\lambda_{\max}$.

For real symmetric $A$ and real $\mathbf x$, all directions in which the quotient is stationary are eigenvectors, not only the extreme ones. By the quotient rule and the gradient $\nabla q_A(\mathbf x)=2A\mathbf x$ (\vcref{1}{sec:differentiability}),
\begin{equation}
\nabla r_A(\mathbf x)=\frac{2A\mathbf x\,\|\mathbf x\|^2-q_A(\mathbf x)\,2\mathbf x}{\|\mathbf x\|^4}=\frac2{\|\mathbf x\|^2}\bigl(A\mathbf x-r_A(\mathbf x)\,\mathbf x\bigr),
\label{eq:rayleigh-gradient}
\end{equation}
so $\nabla r_A(\mathbf x)=\mathbf 0$ exactly when $A\mathbf x=r_A(\mathbf x)\mathbf x$: the critical points of the Rayleigh quotient are the eigenvectors, and the critical values are the eigenvalues. Those for $\lambda_1$ and $\lambda_n$ are the minima and maxima; those for an eigenvalue $\lambda_k$ with $\lambda_1<\lambda_k<\lambda_n$ are saddle points, since moving from $\mathbf q_k$ towards $\mathbf q_1$ lowers the quotient and moving towards $\mathbf q_n$ raises it.

\subsection{Accuracy}

The gradient \eqref{eq:rayleigh-gradient} vanishes at an eigenvector, so to first order the quotient does not change near it. This is the reason for the observation of \cref{tab:power-method} that the Rayleigh quotients of the power method for the symmetric matrix $T_3$ converged twice as fast as the vectors. \Cref{prop:rayleigh-error} proved the bound $|r_A(\mathbf x)-\lambda|\le\|A-\lambda I\|_2\sin^2\angle(\mathbf x,\mathbf x_1)$; with the spectral theorem the error can be written exactly.

\begin{proposition}[The error of the Rayleigh quotient]
\label{prop:rayleigh-quadratic}
Let $A=Q\Lambda Q^{\mathsf H}$ be Hermitian, let $\mathbf q_k$ be one of the eigenvectors, and let $\mathbf x$ be a unit vector with $s\coloneqq\sin\angle(\mathbf x,\mathbf q_k)$. Then, with $c_j\coloneqq\langle\mathbf q_j,\mathbf x\rangle$,
\[
r_A(\mathbf x)-\lambda_k=\sum_{j\ne k}(\lambda_j-\lambda_k)|c_j|^2,\qquad|r_A(\mathbf x)-\lambda_k|\le\max_j|\lambda_j-\lambda_k|\,s^2=\|A-\lambda_kI\|_2\,s^2 .
\]
If $\lambda_k=\lambda_n$, the error $r_A(\mathbf x)-\lambda_n$ is nonpositive, and if $\lambda_k=\lambda_1$, nonnegative.
\end{proposition}

\begin{proof}
Since $\sum_j|c_j|^2=1$, \cref{thm:principal-axes} gives $r_A(\mathbf x)-\lambda_k=\sum_j\lambda_j|c_j|^2-\lambda_k\sum_j|c_j|^2$, and the term $j=k$ cancels. The component of $\mathbf x$ orthogonal to $\mathbf q_k$ is $\sum_{j\ne k}c_j\mathbf q_j$, whose squared length is $s^2=\sum_{j\ne k}|c_j|^2$. The norm $\|A-\lambda_kI\|_2$ equals $\max_j|\lambda_j-\lambda_k|$ by \cref{prop:hermitian-norm} below. If $\lambda_k=\lambda_n$, every term $\lambda_j-\lambda_k$ is $\le0$; if $\lambda_k=\lambda_1$, every term is $\ge0$.
\end{proof}

The spectral theorem also removes the hypothesis of diagonalisability from \cref{thm:power-method}: for a Hermitian matrix with a dominant eigenvalue, the angle decreases like $|\lambda_2/\lambda_1|^k$ and the Rayleigh quotient error like $|\lambda_2/\lambda_1|^{2k}$, in the numbering of that theorem by decreasing modulus. For $T_3$ the ratio is $2/(2+\sqrt2)=2-\sqrt2\approx0.5858$, and its square $(2-\sqrt2)^2=6-4\sqrt2\approx0.3431$ is the ratio observed in the last row of \cref{tab:power-method}. The difference between symmetric and nonsymmetric matrices is shown in \cref{tab:rayleigh-accuracy}. For the eigenvector $\mathbf q$ of $T_3$ for $2+\sqrt2$, normalised, and for the normalised eigenvector $S\mathbf q/\|S\mathbf q\|$ of the matrix $B=ST_3S^{-1}$ of \cref{sec:computing-eigenvalues} for the same eigenvalue, we added $\varepsilon\mathbf w$, where $\mathbf w$ is the unit vector in the direction of the component of $\mathbf e_1$ orthogonal to the eigenvector, and computed the Rayleigh quotient in arbitrary-precision ball arithmetic with python-flint, so that the table shows no rounding errors. For $T_3$ the error is exactly $4\sqrt2\,s^2/3\approx1.886\,s^2$, as the proposition predicts with $c_j$ computed from this $\mathbf w$; for $B$ it is proportional to $s$.

\begin{table}
\sdfloat{%
\begin{sdtable}{colspec={rllll}, row{1-2}={font=\small\bfseries}}
\toprule
 & & \SetCell[c=2]{l} symmetric $T_3$ & & nonsymmetric $B$ \\
\cmidrule[lr]{3-4}\cmidrule[lr]{5-5}
$\varepsilon$ & $s$ & error & error$/s^2$ & error$/s$ \\
\midrule
$10^{-1}$ & $9.95\times10^{-2}$ & $1.87\times10^{-2}$ & $1.886$ & $1.074$ \\
$10^{-2}$ & $1.00\times10^{-2}$ & $1.89\times10^{-4}$ & $1.886$ & $1.284$ \\
$10^{-3}$ & $1.00\times10^{-3}$ & $1.89\times10^{-6}$ & $1.886$ & $1.304$ \\
$10^{-4}$ & $1.00\times10^{-4}$ & $1.89\times10^{-8}$ & $1.886$ & $1.306$ \\
$10^{-6}$ & $1.00\times10^{-6}$ & $1.89\times10^{-12}$ & $1.886$ & $1.307$ \\
$10^{-8}$ & $1.00\times10^{-8}$ & $1.89\times10^{-16}$ & $1.886$ & $1.307$ \\
\bottomrule
\end{sdtable}}{The error $|r(\mathbf x)-(2+\sqrt2)|$ of the Rayleigh quotient of $\mathbf x=\mathbf q+\varepsilon\mathbf w$, for the unit eigenvector $\mathbf q$ of the symmetric matrix $T_3$ and of the nonsymmetric matrix $B$ with the same eigenvalues, and a unit vector $\mathbf w\perp\mathbf q$; $s=\varepsilon/\sqrt{1+\varepsilon^2}$ is the sine of the angle between $\mathbf x$ and $\mathbf q$. Computed in arbitrary precision.\label{tab:rayleigh-accuracy}}
\end{table}

\subsection{The Norm and the Condition Number of a Hermitian Matrix}

\Cref{prop:spectral-radius-norm} showed that $\rho(A)\le\|A\|_2$ for every matrix, and that the inequality can be far from sharp. For Hermitian matrices it is an equality.

\begin{proposition}[The norm of a Hermitian matrix]
\label{prop:hermitian-norm}
For a Hermitian $A\in\mathbb F^{n\times n}$,
\[
\|A\|_2=\rho(A)=\max_j|\lambda_j|=\max(|\lambda_1|,|\lambda_n|),
\]
and if $A$ is invertible, $\|A^{-1}\|_2=1/\min_j|\lambda_j|$. The same holds for every normal matrix, with the complex eigenvalues $\lambda_j$.
\end{proposition}

\begin{proof}
Let $A=Q\Lambda Q^{\mathsf H}$ with unitary $Q$ and diagonal $\Lambda$, as in \cref{thm:spectral} or, for a normal matrix, \cref{thm:normal-unitary}. Since $Q$ and $Q^{\mathsf H}$ preserve lengths, $\|A\mathbf x\|=\|\Lambda Q^{\mathsf H}\mathbf x\|$, and as $\mathbf x$ runs through the unit vectors, so does $\mathbf y=Q^{\mathsf H}\mathbf x$. Hence $\|A\|_2=\|\Lambda\|_2=\max_j|\lambda_j|$, the norm of a diagonal matrix (proof of \cref{thm:bauer-fike}). The inverse is $Q\Lambda^{-1}Q^{\mathsf H}$, with the eigenvalues $1/\lambda_j$. For a Hermitian matrix the largest modulus of an eigenvalue is attained at $\lambda_1$ or at $\lambda_n$.
\end{proof}

For the matrix $A_2$ of \cref{sec:matrix-powers}, which is not normal, $\rho(A_2)=4/5$ and $\|A_2\|_2>2$. For a matrix that is not normal, $\|A\|_2^2$ is the largest eigenvalue of the Hermitian matrix $A^{\mathsf H}A$, by \cref{thm:rayleigh} applied to $r_{A^{\mathsf H}A}(\mathbf x)=\|A\mathbf x\|^2/\|\mathbf x\|^2$; \cref{ch:svd} develops this into the singular value decomposition. The condition number of \cref{sec:householder-qr}, the ratio of the longest to the shortest image $A\mathbf x$ of a unit vector, becomes for an invertible Hermitian matrix
\[
\kappa_2(A)=\|A\|_2\|A^{-1}\|_2=\frac{\max_j|\lambda_j|}{\min_j|\lambda_j|},\qquad\kappa_2(A)=\frac{\lambda_n}{\lambda_1}\quad\text{if }A\succ0 .
\]
For the second difference matrix, with $\lambda_1=4\sin^2(\pi/(2(n+1)))$ and $\lambda_n=4\cos^2(\pi/(2(n+1)))$ (\cref{thm:second-difference-eigen}),
\[
\kappa_2(T_n)=\cot^2\frac{\pi}{2(n+1)}=\frac{4(n+1)^2}{\pi^2}-\frac23+O\bigl((n+1)^{-2}\bigr),
\]
where the expansion follows from $\cot x=1/x-x/3+O(x^3)$ (\vcref{1}{thm:sin-cos-series}), which gives $\cot^2x=1/x^2-2/3+O(x^2)$. \Cref{tab:kappa-tn} compares the exact value, computed with python-flint, with the leading term $4(n+1)^2/\pi^2$ promised in \cref{sec:structured-eigen} and with the value returned by the routine \code{linalg.cond}; the difference between the first two approaches $2/3$.

\begin{table}
\sdfloat{%
\begin{sdtable}{colspec={rllll}, row{1}={font=\small\bfseries}}
\toprule
$n$ & $\kappa_2(T_n)$ & $4(n+1)^2/\pi^2$ & difference & \code{cond}, rel.\ deviation \\
\midrule
$10$ & $48.374150$ & $49.039453$ & $0.665303$ & $1.8\times10^{-15}$ \\
$100$ & $4133.6429$ & $4134.3096$ & $0.666651$ & $3.2\times10^{-16}$ \\
$1000$ & $406095.04$ & $406095.71$ & $0.666667$ & $2.6\times10^{-11}$ \\
\bottomrule
\end{sdtable}}{The condition number $\kappa_2(T_n)=\cot^2(\pi/(2(n+1)))$ of the second difference matrix, its approximation $4(n+1)^2/\pi^2$, their difference, and the relative deviation of the value computed by \code{linalg.cond} in binary64 from the exact one.\label{tab:kappa-tn}}
\end{table}

\section{Min--Max Principles}
\label{sec:min-max}

\subsection{The Courant--Fischer Theorem}

\Cref{thm:rayleigh} characterises the extreme eigenvalues by optimisation. The eigenvalues in between can be characterised in the same way, if the optimisation is over subspaces as well as over vectors. One way is to restrict $\mathbf x$ to be orthogonal to the eigenvectors already found: $\lambda_2=\min\{r_A(\mathbf x):\mathbf x\ne\mathbf 0,\ \mathbf x\perp\mathbf q_1\}$, by the proof of \cref{thm:rayleigh}. This needs the eigenvectors. The following characterisation does not.

\begin{theorem}[Courant--Fischer]
\label{thm:courant-fischer}
Let $A\in\mathbb F^{n\times n}$ be Hermitian, and let $k\in\{1,\ldots,n\}$. Then
\[
\lambda_k(A)=\min_{\dim U=k}\ \max_{\mathbf x\in U\setminus\{\mathbf 0\}}r_A(\mathbf x)=\max_{\dim U=n-k+1}\ \min_{\mathbf x\in U\setminus\{\mathbf 0\}}r_A(\mathbf x),
\]
where $U$ runs through the subspaces of $\mathbb F^n$ of the dimension indicated; the extremes are attained.
\end{theorem}

\begin{proof}
Let $\mathbf q_1,\ldots,\mathbf q_n$ be orthonormal eigenvectors with $A\mathbf q_j=\lambda_j\mathbf q_j$, $\lambda_1\le\cdots\le\lambda_n$. For a nonzero $\mathbf x=\sum_jc_j\mathbf q_j$, $r_A(\mathbf x)=\sum_j\lambda_j|c_j|^2/\sum_j|c_j|^2$ by \cref{thm:principal-axes}; on the span of $\mathbf q_i,\ldots,\mathbf q_l$ this weighted average lies between $\lambda_i$ and $\lambda_l$. The maximum of the continuous function $r_A$ over $U\setminus\{\mathbf 0\}$ exists, since it is the maximum over the unit vectors of $U$, a closed bounded set.

Let $U$ have dimension $k$, and let $V\coloneqq\operatorname{span}\{\mathbf q_k,\ldots,\mathbf q_n\}$, of dimension $n-k+1$. By the dimension formula (\cref{thm:dimension-formula}),
\[
\dim(U\cap V)=\dim U+\dim V-\dim(U+V)\ge k+(n-k+1)-n=1 ,
\]
so $U\cap V$ contains some $\mathbf x\ne\mathbf 0$, for which $r_A(\mathbf x)\ge\lambda_k$ because $\mathbf x\in V$. Hence the maximum over $U$ is at least $\lambda_k$, for every $U$ of dimension $k$. For $U=\operatorname{span}\{\mathbf q_1,\ldots,\mathbf q_k\}$ the maximum is at most $\lambda_k$, and equal to it at $\mathbf q_k$. This proves the first formula. The second is the first applied to $-A$, whose eigenvalues in increasing order are $-\lambda_n\le\cdots\le-\lambda_1$, so that $\lambda_k(A)=-\lambda_{n-k+1}(-A)$ and $r_{-A}=-r_A$.
\end{proof}

For $k=1$ and $k=n$ the theorem reduces to \cref{thm:rayleigh}, since the only subspace of dimension $n$ is $\mathbb F^n$ and the maximum over a line is the value at any of its nonzero vectors. The force of the theorem lies in its consequences, which compare the eigenvalues of different matrices without computing any eigenvectors.

\subsection{Weyl's Inequalities}

\begin{theorem}[Weyl's inequalities]
\label{thm:weyl}
Let $A,E\in\mathbb F^{n\times n}$ be Hermitian. For $k=1,\ldots,n$,
\[
\lambda_k(A)+\lambda_1(E)\le\lambda_k(A+E)\le\lambda_k(A)+\lambda_n(E),
\]
and consequently
\[
|\lambda_k(A+E)-\lambda_k(A)|\le\|E\|_2 .
\]
\end{theorem}

\begin{proof}
For every nonzero $\mathbf x$, $r_{A+E}(\mathbf x)=r_A(\mathbf x)+r_E(\mathbf x)$, and $\lambda_1(E)\le r_E(\mathbf x)\le\lambda_n(E)$ by \cref{thm:rayleigh}. Hence, for each subspace $U$, the maximum of $r_{A+E}$ over $U$ lies between the maximum of $r_A$ plus $\lambda_1(E)$ and the maximum of $r_A$ plus $\lambda_n(E)$, and taking the minimum over all $U$ of dimension $k$ gives the two inequalities by \cref{thm:courant-fischer}. Since $\|E\|_2=\max(|\lambda_1(E)|,|\lambda_n(E)|)$ by \cref{prop:hermitian-norm}, both $\lambda_n(E)$ and $-\lambda_1(E)$ are at most $\|E\|_2$.
\end{proof}

The second inequality says that the eigenvalues of Hermitian matrices are perfectly conditioned: a Hermitian perturbation of norm $\|E\|_2$ moves each eigenvalue, the $k$th smallest to the $k$th smallest, by at most $\|E\|_2$. It is sharper than the Bauer--Fike theorem with $\kappa_2(Q)=1$, which only says that every eigenvalue of $A+E$ is within $\|E\|_2$ of \emph{some} eigenvalue of $A$, and it is the opposite of what \cref{sec:computing-eigenvalues} found for the Jordan block $J_n(0)$, whose eigenvalue $0$ moves by $\varepsilon^{1/n}$ under a perturbation of size $\varepsilon$. It explains the experiment of \cref{tab:tn-eigen} as far as the eigenvalues of the matrix are concerned: if a computed set of eigenvalues is the exact set of eigenvalues of a Hermitian matrix $T_n+E$ with $\|E\|_2$ of the order of $u\|T_n\|_2$, as Volume~IX proves for the standard methods~\cite[Sect.~8.3]{golub2013}, its absolute errors are of that order, and small eigenvalues have correspondingly large relative errors. \Cref{sec:symmetric-computation} tests the theorem.

Two special cases are often used. If $E\succeq0$, then $\lambda_1(E)\ge0$, and every eigenvalue increases or stays the same.

\begin{corollary}[Monotonicity]
\label{cor:monotone}
If $A$ and $B$ are Hermitian and $A\preceq B$, then $\lambda_k(A)\le\lambda_k(B)$ for $k=1,\ldots,n$. If $A\prec B$, the inequalities are strict.
\end{corollary}

\begin{proof}
Apply \cref{thm:weyl} to $A$ and $E=B-A$, whose smallest eigenvalue is nonnegative, or positive if $A\prec B$, by \cref{thm:principal-axes}.
\end{proof}

\subsection{Interlacing}

Deleting a row and the corresponding column of a Hermitian matrix leaves a Hermitian matrix whose eigenvalues interlace with those of the original one.

\begin{theorem}[Cauchy's interlacing theorem]
\label{thm:interlacing}
Let $A\in\mathbb F^{n\times n}$ be Hermitian, and let $B=A_S$ be a principal submatrix of order $m$, formed by the rows and the columns with indices in $S$, $|S|=m$. Then
\[
\lambda_k(A)\le\lambda_k(B)\le\lambda_{k+n-m}(A)\qquad(k=1,\ldots,m).
\]
In particular, for $m=n-1$,
\[
\lambda_1(A)\le\lambda_1(B)\le\lambda_2(A)\le\lambda_2(B)\le\cdots\le\lambda_{n-1}(B)\le\lambda_n(A).
\]
\end{theorem}

\begin{proof}
Let $E_S\in\mathbb F^{n\times m}$ have the columns $\mathbf e_i$, $i\in S$, so that $B=E_S^{\mathsf H}AE_S$, and let $V_S\coloneqq\mathcal C(E_S)=\operatorname{span}\{\mathbf e_i:i\in S\}$. The map $\mathbf y\mapsto E_S\mathbf y$ preserves lengths and maps the subspaces of $\mathbb F^m$ of dimension $k$ bijectively onto the subspaces of $V_S$ of dimension $k$, and $r_B(\mathbf y)=r_A(E_S\mathbf y)$. By \cref{thm:courant-fischer}, therefore,
\[
\lambda_k(B)=\min_{\substack{U\subseteq V_S\\\dim U=k}}\ \max_{\mathbf x\in U\setminus\{\mathbf 0\}}r_A(\mathbf x)\ge\min_{\dim U=k}\ \max_{\mathbf x\in U\setminus\{\mathbf 0\}}r_A(\mathbf x)=\lambda_k(A),
\]
since the minimum on the left is over fewer subspaces. In the same way the second formula of the theorem, with subspaces of dimension $m-k+1$ of $V_S$, gives
\[
\lambda_k(B)=\max_{\substack{U\subseteq V_S\\\dim U=m-k+1}}\ \min_{\mathbf x\in U\setminus\{\mathbf 0\}}r_A(\mathbf x)\le\max_{\dim U=m-k+1}\ \min_{\mathbf x\in U\setminus\{\mathbf 0\}}r_A(\mathbf x)=\lambda_{n-(m-k+1)+1}(A),
\]
which is $\lambda_{k+n-m}(A)$.
\end{proof}

\Cref{fig:interlacing} shows the theorem for the second difference matrices: $T_{k}$ is the leading principal submatrix of order $k$ of $T_{k+1}$, so the eigenvalues of $T_1,T_2,\ldots,T_6$ interlace from each row to the next. The interlacing theorem also gives a second proof that leading principal minors of a positive definite matrix are positive, since all eigenvalues of $A_k$ are at least $\lambda_1(A)>0$, and it shows that a Hermitian matrix with a negative principal minor of order $2$, or a negative diagonal entry, has a negative eigenvalue: $\lambda_1(A)\le\lambda_1(B)<0$ for that principal submatrix $B$ if $\det B<0$ or $B=[a_{ii}]$ with $a_{ii}<0$.

\begin{figure}
\sdfloat{%
\begin{tikzpicture}
\begin{axis}[sdplot, width=0.62\textwidth, height=0.44\textwidth, xmin=0, xmax=4, ymin=0.4, ymax=6.6,
  xtick={0,1,2,3,4}, ytick={1,2,3,4,5,6}, yticklabels={$T_1$,$T_2$,$T_3$,$T_4$,$T_5$,$T_6$}, xlabel={eigenvalue},
  axis y line*=left, y axis line style={draw=none}, ytick style={draw=none}]
\foreach \k in {1,...,6} {
\addplot[black!25, line width=0.4pt, domain=0:4, samples=2] {\k};
}
\addplot[black!70, dashed, line width=0.45pt] coordinates {(2,1) (1,2) (0.585786,3) (0.381966,4) (0.267949,5) (0.198062,6)};
\addplot[black!70, dashed, line width=0.45pt] coordinates {(2,1) (3,2) (3.414214,3) (3.618034,4) (3.732051,5) (3.801938,6)};
\addplot[black, only marks, mark=*, mark size=1.5pt, samples at={1}] ({2-2*cos(deg(x*pi/2))}, 1);
\addplot[black, only marks, mark=*, mark size=1.5pt, samples at={1,2}] ({2-2*cos(deg(x*pi/3))}, 2);
\addplot[black, only marks, mark=*, mark size=1.5pt, samples at={1,2,3}] ({2-2*cos(deg(x*pi/4))}, 3);
\addplot[black, only marks, mark=*, mark size=1.5pt, samples at={1,...,4}] ({2-2*cos(deg(x*pi/5))}, 4);
\addplot[black, only marks, mark=*, mark size=1.5pt, samples at={1,...,5}] ({2-2*cos(deg(x*pi/6))}, 5);
\addplot[black, only marks, mark=*, mark size=1.5pt, samples at={1,...,6}] ({2-2*cos(deg(x*pi/7))}, 6);
\end{axis}
\end{tikzpicture}}{The eigenvalues $2-2\cos(j\pi/(k+1))$, $j=1,\ldots,k$, of the second difference matrices $T_k$ for $k=1,\ldots,6$, one row for each $k$. Between two neighbouring eigenvalues of $T_{k+1}$ lies exactly one eigenvalue of $T_k$, its leading principal submatrix (\cref{thm:interlacing}); the dashed lines join the smallest and the largest eigenvalues, which move outwards as $k$ grows.\label{fig:interlacing}}
\end{figure}

A perturbation of rank one interlaces in the same way.

\begin{proposition}[Rank-one perturbations]
\label{prop:rank-one-interlacing}
Let $A\in\mathbb F^{n\times n}$ be Hermitian, $\mathbf v\in\mathbb F^n$, and $B\coloneqq A+\mathbf v\mathbf v^{\mathsf H}$. Then
\[
\lambda_k(A)\le\lambda_k(B)\le\lambda_{k+1}(A)\quad(k=1,\ldots,n-1),\qquad\lambda_n(A)\le\lambda_n(B),
\]
and $\sum_k\bigl(\lambda_k(B)-\lambda_k(A)\bigr)=\|\mathbf v\|^2$.
\end{proposition}

\begin{proof}
The matrix $\mathbf v\mathbf v^{\mathsf H}$ is positive semidefinite, so $\lambda_k(A)\le\lambda_k(B)$ by \cref{cor:monotone}. For the upper bound, let $k\le n-1$, $V\coloneqq\operatorname{span}\{\mathbf q_1,\ldots,\mathbf q_{k+1}\}$ for orthonormal eigenvectors of $A$, and $U\coloneqq V\cap\{\mathbf v\}^\perp$, which has dimension at least $k$ by the dimension formula; let $U'\subseteq U$ have dimension exactly $k$. On $U'$, $r_B(\mathbf x)=r_A(\mathbf x)+|\langle\mathbf v,\mathbf x\rangle|^2/\|\mathbf x\|^2=r_A(\mathbf x)\le\lambda_{k+1}(A)$, since $\mathbf x\in V$. By \cref{thm:courant-fischer}, $\lambda_k(B)\le\max_{U'}r_B\le\lambda_{k+1}(A)$. The sum is $\operatorname{tr}B-\operatorname{tr}A=\operatorname{tr}(\mathbf v\mathbf v^{\mathsf H})=\|\mathbf v\|^2$, by \cref{thm:trace-det-eigen}.
\end{proof}

A positive rank-one change thus pushes every eigenvalue up, but no eigenvalue past the next one of the original matrix, and the total increase is $\|\mathbf v\|^2$. Such changes occur when an observation is added to a least squares problem, $X^{\mathsf T}X\mapsto X^{\mathsf T}X+\mathbf x\mathbf x^{\mathsf T}$ (\cref{thm:sherman-morrison}).

\section{Inertia and Indefinite Forms}
\label{sec:inertia}

\subsection{Sylvester's Law of Inertia}

A congruence $A\mapsto S^{\mathsf H}AS$ changes the eigenvalues, as the example $S^{\mathsf T}IS=\operatorname{diag}(4,1)$ of \cref{sec:quadratic-forms} shows. It cannot change their signs.

\begin{definition}[Inertia]
\label{def:inertia}
The \term{inertia} of a Hermitian matrix $A$ is the triple $(n_+,n_-,n_0)$ of the numbers of its positive, negative and zero eigenvalues, counted with multiplicity.
\end{definition}

\begin{theorem}[Sylvester's law of inertia]
\label{thm:inertia}
Congruent Hermitian matrices have the same inertia. More precisely, for a Hermitian $A\in\mathbb F^{n\times n}$, $n_+(A)$ is the largest dimension of a subspace $U$ of $\mathbb F^n$ with $\mathbf x^{\mathsf H}A\mathbf x>0$ for all $\mathbf x\in U\setminus\{\mathbf 0\}$, $n_-(A)$ the largest dimension of one with $\mathbf x^{\mathsf H}A\mathbf x<0$, and $n_0(A)=\dim\mathcal N(A)$. Every Hermitian matrix is congruent to
\[
\operatorname{diag}(I_{n_+},-I_{n_-},0_{n_0}),
\]
and two Hermitian matrices of the same order are congruent if and only if they have the same inertia.
\end{theorem}

\begin{proof}
Let $A=Q\Lambda Q^{\mathsf H}$. On the span $U_+$ of the eigenvectors for the positive eigenvalues, the form is $\sum\lambda_j|c_j|^2>0$ away from $\mathbf 0$ (\cref{thm:principal-axes}), and $\dim U_+=n_+(A)$. If $U$ is any subspace on which the form is positive away from $\mathbf 0$, then $U$ meets the span $W$ of the eigenvectors for the eigenvalues $\le0$, of dimension $n-n_+(A)$, only in $\mathbf 0$, since the form is $\le0$ on $W$; by the dimension formula, $\dim U+(n-n_+(A))-\dim(U\cap W)\le n$, so $\dim U\le n_+(A)$. This proves the characterisation of $n_+$, and that of $n_-$ follows with $-A$. The number of zero eigenvalues is the geometric multiplicity of $0$, which is $\dim\mathcal N(A)$.

If $B=S^{\mathsf H}AS$ with an invertible $S$, then $\mathbf y^{\mathsf H}B\mathbf y=(S\mathbf y)^{\mathsf H}A(S\mathbf y)$, so $S$ maps every subspace on which $q_B$ is positive onto a subspace of the same dimension on which $q_A$ is positive, and $S^{-1}$ does the converse; hence $n_+(B)=n_+(A)$, likewise $n_-(B)=n_-(A)$, and $n_0=n-n_+-n_-$. For the normal form, reorder the columns of $Q$ so that the positive eigenvalues come first, then the negative ones, and let $D$ be the diagonal matrix with the entries $|\lambda_j|^{-1/2}$ for $\lambda_j\ne0$ and $1$ for $\lambda_j=0$; then $(QD)^{\mathsf H}A(QD)=D\Lambda D=\operatorname{diag}(I_{n_+},-I_{n_-},0)$. Two matrices with the same inertia are congruent to the same normal form and hence to each other, because congruence is an equivalence relation.
\end{proof}

The inertia is thus the complete invariant of a Hermitian form under changes of variables: up to a change of coordinates there are only the forms $|y_1|^2+\cdots+|y_p|^2-|y_{p+1}|^2-\cdots-|y_{p+q}|^2$. The numbers $n_+-n_-$ and $n_++n_-=\operatorname{rank}A$ are called the signature and the rank of the form.

\subsection{Inertia from Elimination}

Since $A=LDL^{\mathsf H}$ is a congruence when $L$ is invertible, the law of inertia lets us read the signs of the eigenvalues from the pivots, without computing any eigenvalue.

\begin{proposition}[Inertia from pivots]
\label{prop:inertia-ldl}
If a Hermitian matrix has a factorisation $A=LDL^{\mathsf H}$ with an invertible $L$ and a real diagonal $D$, then the inertia of $A$ is the triple of the numbers of positive, negative and zero diagonal entries of $D$. In particular, if Gaussian elimination without exchanges can be carried out on $A$ with nonzero pivots, $n_-(A)$ is the number of negative pivots.
\end{proposition}

\begin{proof}
$D=L^{-1}AL^{-\mathsf H}$ is congruent to $A$, and the eigenvalues of a diagonal matrix are its diagonal entries. The second statement is the case of \cref{prop:ldu} in its Hermitian form, as in the proof of \cref{thm:pd-criteria}.
\end{proof}

This is the classification at a critical point of \vcref{1}{sec:extrema-n} in its final form. If the Hessian $\nabla^2f(\mathbf c)$ is invertible, its inertia is $(n_+,n_-,0)$, and $n_-$, the number of independent directions in which $f$ decreases to second order, is called the \term{index} of the critical point: $0$ at a minimum, $n$ at a maximum, and in between at the various kinds of saddle points. It can be computed by elimination on the Hessian, for instance by completing the square. For $f(x,y)=x^3-3xy+y^3$ at the saddle point $(0,0)$, the Hessian has the rows $(0,-3)$ and $(-3,0)$, a zero first pivot; but its determinant $-9$ is the product of the two eigenvalues, which therefore have opposite signs, and the inertia is $(1,1,0)$.

\subsection{Counting Eigenvalues}

The eigenvalues of $A-\sigma I$ are the numbers $\lambda_j-\sigma$, so $n_-(A-\sigma I)$ is the number of eigenvalues of $A$ less than $\sigma$.

\begin{proposition}[Counting eigenvalues below a shift]
\label{prop:eigen-count}
Let $A$ be Hermitian and $\sigma$ real. If elimination without exchanges can be carried out on $A-\sigma I$ with nonzero pivots $d_1,\ldots,d_n$, then the number of eigenvalues of $A$ less than $\sigma$ is the number of negative $d_k$. For a real symmetric tridiagonal matrix with the diagonal entries $a_1,\ldots,a_n$ and the entries $b_1,\ldots,b_{n-1}$ next to the diagonal, the pivots are
\[
d_1=a_1-\sigma,\qquad d_k=a_k-\sigma-\frac{b_{k-1}^2}{d_{k-1}}\quad(k=2,\ldots,n),
\]
and they are nonzero for all $\sigma$ except the finitely many that are eigenvalues of one of the leading principal submatrices of $A$.
\end{proposition}

\begin{proof}
The first statement is \cref{prop:inertia-ldl} applied to $A-\sigma I$. In elimination on a tridiagonal matrix, step $k-1$ subtracts $b_{k-1}/d_{k-1}$ times row $k-1$ from row $k$, which changes only the diagonal entry of row $k$, by $-b_{k-1}^2/d_{k-1}$, as for the Thomas algorithm (\cref{alg:thomas}). By \cref{prop:pivots-minors}, $d_k=\det(A_k-\sigma I)/\det(A_{k-1}-\sigma I)$ as long as the previous pivots are nonzero, so $d_k\ne0$ for all $k$ if $\sigma$ is not an eigenvalue of any $A_k$.
\end{proof}

\begin{algorithm}
\begin{algorithmic}[1]
\Require the diagonal $a_1,\ldots,a_n$ and the entries $b_1,\ldots,b_{n-1}$ next to it of a real symmetric tridiagonal matrix $A$, and a real $\sigma$ for which all pivots are nonzero
\Ensure the number of eigenvalues of $A$ less than $\sigma$
\Name{CountBelow}{a,b,\sigma}
\State $d\gets a_1-\sigma$
\State $c\gets1$ if $d<0$, else $c\gets0$
\For{$k=2,\ldots,n$}
  \State $d\gets a_k-\sigma-b_{k-1}^2/d$ \Comment{the $k$th pivot of $A-\sigma I$}
  \If{$d<0$} $c\gets c+1$ \EndIf
\EndFor
\State \Return $c$
\end{algorithmic}
\caption{Counting the eigenvalues of a symmetric tridiagonal matrix below $\sigma$ by the signs of the pivots of $A-\sigma I$ (\cref{prop:eigen-count}). It costs about $4n$ operations. Since the count is a nondecreasing function of $\sigma$, bisection on $\sigma$ locates each eigenvalue to any accuracy.}
\label{alg:eigen-count}
\end{algorithm}

For $T_5$ and $\sigma=3/2$ the pivots are
\begin{gather*}
d_1=\frac12,\qquad d_2=\frac12-2=-\frac32,\qquad d_3=\frac12+\frac23=\frac76,\\
d_4=\frac12-\frac67=-\frac5{14},\qquad d_5=\frac12+\frac{14}5=\frac{33}{10},
\end{gather*}
two of them negative; indeed $T_5$ has the eigenvalues $2-\sqrt3$, $1$, $2$, $3$ and $2+\sqrt3$ (\cref{thm:second-difference-eigen}), two of which are less than $3/2$. The product of the pivots is $33/32=\det(T_5-3I/2)$, which is also the product of the five numbers $\lambda_j-3/2$. Counting is the basis of the \term{bisection method} for eigenvalues of symmetric tridiagonal matrices: if $c(\sigma)$ denotes the count, then $\lambda_k$ lies in an interval $[\alpha,\beta]$ with $c(\alpha)<k\le c(\beta)$, and halving the interval repeatedly by evaluating $c$ at its midpoint locates $\lambda_k$ to any desired accuracy. The pivots are related to the leading principal minors $\det(A_k-\sigma I)$, a Sturm sequence, through \cref{prop:pivots-minors}, and the method goes back to Givens~\cite{givens1954}; Volume~IX treats it with the reduction of a symmetric matrix to tridiagonal form, by which it becomes a method for all symmetric matrices~\cite[Sect.~8.4]{golub2013}.

\subsection{Indefinite Forms}

Indefinite forms are not defects; some of the most important forms are indefinite. In the special theory of relativity, events are points $(t,x,y,z)$ of space-time, and the \term{Minkowski form}
\[
\langle\mathbf u,\mathbf v\rangle_{\mathrm M}=c^2u_tv_t-u_xv_x-u_yv_y-u_zv_z ,
\]
with the speed of light $c$, takes the place of the inner product. Its matrix, the diagonal matrix with the entries $c^2$, $-1$, $-1$ and $-1$, has the inertia $(1,3,0)$, and by \cref{thm:inertia} no change of coordinates turns it into a definite form; the vectors with $\langle\mathbf u,\mathbf u\rangle_{\mathrm M}>0$, $<0$ and $=0$ are separated by the light cone $\langle\mathbf u,\mathbf u\rangle_{\mathrm M}=0$, and the linear maps that preserve the form, the Lorentz transformations, are the counterpart of the orthogonal maps. This is the indefinite form of \cref{sec:inner-product-spaces}, which fails to be an inner product because positivity fails, and the law of inertia classifies all such forms by their signs.

In optimisation, the minimisation of a quadratic function $\langle\mathbf x,H\mathbf x\rangle/2-\langle\mathbf c,\mathbf x\rangle$ subject to linear constraints $B\mathbf x=\mathbf d$, with $H\in\mathbb R^{n\times n}$ symmetric positive definite and $B\in\mathbb R^{m\times n}$ of rank $m$, leads by the method of Lagrange multipliers to the linear system
\[
K\begin{bmatrix}\mathbf x\\\boldsymbol\lambda\end{bmatrix}=\begin{bmatrix}\mathbf c\\\mathbf d\end{bmatrix},\qquad K\coloneqq\begin{bmatrix}H&B^{\mathsf T}\\B&0\end{bmatrix},
\]
whose symmetric matrix is always indefinite. Indeed, $K$ is congruent to the block diagonal matrix with the blocks $H$ and $-BH^{-1}B^{\mathsf T}$, by the factorisation in the proof of \cref{thm:schur-pd}, and $BH^{-1}B^{\mathsf T}$ is positive definite by \cref{prop:pd-properties}\,(iii),(iv), since $B^{\mathsf T}$ has independent columns; so the inertia of $K$ is $(n,m,0)$, and $K$ is invertible. Such matrices, which Volume~XIII meets in the Karush--Kuhn--Tucker conditions, cannot be factorised by Cholesky's method; symmetric indefinite matrices are factorised as $PAP^{\mathsf T}=LDL^{\mathsf T}$ with symmetric row and column exchanges and a block diagonal $D$ with blocks of order $1$ and $2$, by the method of Bunch and Kaufman~\cite{bunchkaufman1977}, which Volume~IX describes with the routine \code{linalg.ldl}.

\subsection{The Generalised Eigenvalue Problem}

Two quadratic forms, one of them positive definite, can be diagonalised by one change of variables.

\begin{theorem}[Simultaneous diagonalisation by congruence]
\label{thm:generalised-eigen}
Let $A,B\in\mathbb F^{n\times n}$ be Hermitian, with $B$ positive definite. There is an invertible $S\in\mathbb F^{n\times n}$, real if $A$ and $B$ are real, with
\[
S^{\mathsf H}BS=I,\qquad S^{\mathsf H}AS=\Lambda=\operatorname{diag}(\lambda_1,\ldots,\lambda_n),
\]
$\lambda_j$ real. The columns $\mathbf s_j$ of $S$ satisfy $A\mathbf s_j=\lambda_jB\mathbf s_j$ and $\mathbf s_i^{\mathsf H}B\mathbf s_j=\delta_{ij}$, and the $\lambda_j$ are the eigenvalues of $B^{-1}A$, the roots of $\det(A-\lambda B)=0$.
\end{theorem}

\begin{proof}
Let $B=R^{\mathsf H}R$ be the Cholesky factorisation, and let $C\coloneqq R^{-\mathsf H}AR^{-1}$, which is Hermitian. By \cref{thm:spectral}, $C=Q\Lambda Q^{\mathsf H}$ with a unitary $Q$, and $S\coloneqq R^{-1}Q$ satisfies $S^{\mathsf H}BS=Q^{\mathsf H}R^{-\mathsf H}R^{\mathsf H}RR^{-1}Q=I$ and $S^{\mathsf H}AS=Q^{\mathsf H}CQ=\Lambda$. From $S^{\mathsf H}BS=I$, $S^{-\mathsf H}=BS$, and $S^{\mathsf H}AS=\Lambda$ gives $AS=S^{-\mathsf H}\Lambda=BS\Lambda$, whose $j$th column is $A\mathbf s_j=\lambda_jB\mathbf s_j$. Finally $B^{-1}A=S\Lambda S^{-1}$, and $\det(A-\lambda B)=\det(S^{-\mathsf H}(\Lambda-\lambda I)S^{-1})=|\det S|^{-2}\prod_j(\lambda_j-\lambda)$.
\end{proof}

The matrix $B^{-1}A$ is not Hermitian in general, but the theorem shows that it behaves like one: its eigenvalues are real, and its eigenvectors are orthonormal in the inner product $\langle\cdot,\cdot\rangle_B$ of \cref{cor:inner-products-pd}, for which $B^{-1}A$ is self-adjoint, since $\langle\mathbf x,B^{-1}A\mathbf y\rangle_B=\mathbf x^{\mathsf H}A\mathbf y=\langle B^{-1}A\mathbf x,\mathbf y\rangle_B$. The proof is also how the problem is solved in computation: one Cholesky factorisation reduces it to an ordinary Hermitian eigenvalue problem, and the routines \code{linalg.eigh} and \code{linalg.eigvalsh} of the library accept the second matrix $B$ for this purpose. Unlike \cref{thm:simultaneous-unitary}, the theorem needs no commutativity, because $S$ is not unitary.

The theorem describes small vibrations. For the chain of masses of \cref{sec:quadratic-forms}, or any mechanical system with the kinetic energy $\langle\dot{\mathbf x},M\dot{\mathbf x}\rangle/2$ and the potential energy $\langle\mathbf x,K\mathbf x\rangle/2$ with symmetric positive definite $M$ and $K$, the equations of motion are $M\ddot{\mathbf x}+K\mathbf x=\mathbf 0$. With $\mathbf x=S\mathbf z$ for the matrix $S$ of the theorem applied to $K$ and $M$, multiplying by $S^{\mathsf T}$ gives $\ddot{\mathbf z}+\Lambda\mathbf z=\mathbf 0$, that is $\ddot z_j+\lambda_jz_j=0$ for each $j$, with $\lambda_j=\mathbf s_j^{\mathsf T}K\mathbf s_j>0$: the system decouples into $n$ independent oscillators with the frequencies $\omega_j=\sqrt{\lambda_j}$, and the columns $\mathbf s_j$ are the \term{normal modes}. For the chain with $M=mI$ and $K=kT_n$, the modes are the sine vectors of \cref{thm:second-difference-eigen} and the frequencies are $\sqrt{(k/m)(2-2\cos(j\pi/(n+1)))}$. \Cref{ch:linear-ode-systems} solves such systems of differential equations in general.

\section{Gaussian Integrals and Other Uses}
\label{sec:gaussian-uses}

\subsection{The Gaussian Integral in \texorpdfstring{$n$}{n} Variables}

\Vcref{1}{thm:gaussian-quadratic} stated that $\int_{\mathbb R^n}e^{-\langle\mathbf x,A\mathbf x\rangle}d\mathbf x=\pi^{n/2}/\sqrt{\det A}$ for every symmetric positive definite $A$, and proved it for $n=2$ by completing the square. The spectral theorem gives the proof for every $n$. It needs the change of variables formula only for affine maps, and that case follows from \vcref{1}{thm:linear-volume} alone.

\begin{lemma}[Affine change of variables]
\label{lem:affine-substitution}
Let $A\in\mathbb R^{n\times n}$ be invertible, $\mathbf c\in\mathbb R^n$ and $\mathbf g(\mathbf u)\coloneqq A\mathbf u+\mathbf c$. For every closed, bounded, Jordan measurable $K\subseteq\mathbb R^n$ and every function $f$ continuous on $\mathbf g(K)$,
\[
\int_{\mathbf g(K)}f(\mathbf x)\,d\mathbf x=|\det A|\int_Kf\bigl(\mathbf g(\mathbf u)\bigr)\,d\mathbf u .
\]
\end{lemma}

\begin{proof}
By \vcref{1}{prop:scaling} with $c=1$, translations preserve integrals, so we may assume $\mathbf c=\mathbf 0$. The set $A(K)$ is closed, bounded and Jordan measurable by \vcref{1}{thm:linear-volume}, and both integrands are continuous, so both integrals exist (\vcref{1}{thm:multiple-summary}). Let $\varepsilon>0$. The function $h\coloneqq f\circ\mathbf g$ is uniformly continuous on the closed bounded set $K$ (\vcref{1}{lem:heine-cantor-rn}); choose $\delta>0$ such that $|h(\mathbf u)-h(\mathbf v)|<\varepsilon$ for $\mathbf u,\mathbf v\in K$ at distance less than $\delta$, and let $M$ be a bound for $|f|$ on $A(K)$. Cover a box containing $K$ by a grid of closed cubes of side $\eta$ and diameter $\sqrt n\,\eta<\delta$, let $Q_1,\ldots,Q_N$ be the cubes contained in $K$, and let $Z\coloneqq K\setminus(Q_1\cup\cdots\cup Q_N)$. A point of $Z$ lies in a cube of the grid that is not contained in $K$ but meets $K$, and such a cube meets $\partial K$, since a segment joining a point of $K$ to a point outside $K$ within the cube contains a boundary point. Cover $\partial K$, which has content zero, by finitely many boxes of total volume less than $\varepsilon/2$; the cubes of the grid that meet $\partial K$ lie in these boxes enlarged by $\eta$ on every side, so for small $\eta$ their total volume is less than $\varepsilon$, as in the proof of \vcref{1}{lem:lipschitz-content}, and $\operatorname{vol}(Z)<\varepsilon$. The boundary of $Z$ lies in $\partial K$ and the faces of the $Q_i$, so $Z$ is Jordan measurable (\vcref{1}{thm:multiple-summary}). By \vcref{1}{thm:linear-volume}, the sets $A(Q_i)$ and $A(Z)$ are Jordan measurable, with $\operatorname{vol}(A(Q_i))=|\det A|\operatorname{vol}(Q_i)$ and $\operatorname{vol}(A(Z))<|\det A|\,\varepsilon$; they make up $A(K)$, and since $A$ is injective, two of them meet only in the image of a common face of two cubes, which has content zero by \vcref{1}{thm:multiple-summary,lem:lipschitz-content}. So the integrals over $A(K)$ and over $K$ split into the sums of the integrals over these pieces (\vcref{1}{thm:integral-over-region}\,(ii), in its $n$-dimensional form). On $Q_i$, $h$ lies within $\varepsilon$ of its value $h_i$ at the centre of $Q_i$, so $f$ lies within $\varepsilon$ of $h_i$ on $A(Q_i)$, and
\begin{gather*}
\Bigl|\int_{A(Q_i)}f-h_i\,|\det A|\operatorname{vol}(Q_i)\Bigr|\le\varepsilon\,|\det A|\operatorname{vol}(Q_i),\\
\Bigl|\int_{Q_i}h-h_i\operatorname{vol}(Q_i)\Bigr|\le\varepsilon\operatorname{vol}(Q_i).
\end{gather*}
The pieces over $A(Z)$ and $Z$ are at most $M|\det A|\varepsilon$ and $M\varepsilon$ in modulus, by \vcref{1}{thm:integral-over-region}\,(i). Adding over $i$ gives $\bigl|\int_{A(K)}f-|\det A|\int_Kh\bigr|\le|\det A|\,\varepsilon\,(2\operatorname{vol}(K)+2M)$, and $\varepsilon$ was arbitrary.
\end{proof}

\begin{theorem}[Gaussian integrals of quadratic forms]
\label{thm:gaussian-general}
Let $A\in\mathbb R^{n\times n}$ be symmetric and positive definite, and $\mathbf b\in\mathbb R^n$. Then
\[
\int_{\mathbb R^n}e^{-\langle\mathbf x,A\mathbf x\rangle}\,d\mathbf x=\frac{\pi^{n/2}}{\sqrt{\det A}},\qquad\int_{\mathbb R^n}e^{-\langle\mathbf x,A\mathbf x\rangle+\langle\mathbf b,\mathbf x\rangle}\,d\mathbf x=\frac{\pi^{n/2}}{\sqrt{\det A}}\,e^{\langle\mathbf b,A^{-1}\mathbf b\rangle/4}.
\]
\end{theorem}

\begin{proof}
Let $A=Q\Lambda Q^{\mathsf T}$ with a real orthogonal $Q$ (\cref{thm:spectral}). The orthogonal map $\mathbf u\mapsto Q\mathbf u$ preserves lengths and maps the closed ball $\bar B(\mathbf 0,j)$ onto itself, and $\langle Q\mathbf u,AQ\mathbf u\rangle=\sum_i\lambda_iu_i^2$ by \cref{thm:principal-axes}. By \cref{lem:affine-substitution}, with $|\det Q|=1$,
\[
\int_{\bar B(\mathbf 0,j)}e^{-\langle\mathbf x,A\mathbf x\rangle}\,d\mathbf x=\int_{\bar B(\mathbf 0,j)}e^{-(\lambda_1u_1^2+\cdots+\lambda_nu_n^2)}\,d\mathbf u .
\]
The balls $\bar B(\mathbf 0,j)$ form an exhaustion of $\mathbb R^n$, and the integrands are positive and continuous; by \vcref{1}{thm:exhaustion} the limit as $j\to\infty$ on the right is the improper integral computed over cubes in \vcref{1}{prop:gaussian-n}, which is $\prod_i\sqrt{\pi/\lambda_i}=\pi^{n/2}/\sqrt{\lambda_1\cdots\lambda_n}$, and $\lambda_1\cdots\lambda_n=\det A$. For the second integral, \cref{prop:quadratic-centre} with $-A$ in place of $A$ and $\mathbf x_0\coloneqq A^{-1}\mathbf b/2$ gives
\[
-\langle\mathbf x,A\mathbf x\rangle+\langle\mathbf b,\mathbf x\rangle=-\langle\mathbf x-\mathbf x_0,A(\mathbf x-\mathbf x_0)\rangle+\tfrac14\langle\mathbf b,A^{-1}\mathbf b\rangle ,
\]
and the translation $\mathbf x=\mathbf x_0+\mathbf u$ maps the exhaustion by the balls $\bar B(\mathbf 0,j)$ to the exhaustion by the balls $\bar B(\mathbf x_0,j)$; by \cref{lem:affine-substitution} and \vcref{1}{thm:exhaustion} the integral is $e^{\langle\mathbf b,A^{-1}\mathbf b\rangle/4}$ times the first one.
\end{proof}

This completes the proof of \vcref{1}{thm:gaussian-quadratic}, and with it the normalisation of the multivariate normal density stated after that theorem, for every $n$. The second formula, read as a function of $\mathbf b$, is the moment generating function of a normal distribution, from which Volume~IV derives its moments. The factorisation $A=LDL^{\mathsf T}$ would serve equally well in the proof, with the map $\mathbf u\mapsto L^{-\mathsf T}\mathbf u$ of determinant $1$ in place of $Q$; it does not map balls onto balls, but the images of the cubes $[-j,j]^n$ form an exhaustion as well.

\subsection{Covariance Matrices}

For random variables $X_1,\ldots,X_n$ with finite variances, the covariance matrix $\Sigma$ has the entries $\sigma_{ij}=\operatorname{cov}(X_i,X_j)$, and for every $\mathbf w\in\mathbb R^n$ the variance of the combination $\langle\mathbf w,\mathbf X\rangle$ is $\langle\mathbf w,\Sigma\mathbf w\rangle$ (Volume~IV). Variances are nonnegative, so every covariance matrix is positive semidefinite, and it is singular exactly when some nontrivial combination of the variables has variance $0$, that is, is constant with probability $1$. The same holds for data. If the $N$ rows of $X\in\mathbb R^{N\times n}$ are observations of $n$ variables, centred by subtracting the mean of each column, the sample covariance matrix $S=X^{\mathsf T}X/(N-1)$ is a Gram matrix, positive semidefinite by \cref{thm:psd-criteria}\,(iii) and positive definite exactly when the columns of $X$ are independent; for $N\le n$ it is always singular, since $\operatorname{rank}X\le N-1$ after centring.

The spectral theorem gives the covariance matrix a geometric meaning. For a unit vector $\mathbf w$, $\langle\mathbf w,S\mathbf w\rangle$ is the sample variance of the projections $\langle\mathbf w,\mathbf x_i\rangle$ of the observations onto the line spanned by $\mathbf w$, and by \cref{thm:rayleigh} the direction of largest variance is the eigenvector of $S$ for $\lambda_{\max}$, and the largest variance is $\lambda_{\max}$. By \cref{thm:courant-fischer} the eigenvectors for the next eigenvalues give the directions of largest variance orthogonal to the previous ones. These are the \term{principal components} of the data, which \cref{ch:svd} computes from the singular value decomposition of $X$ and Volume~XVI uses in statistical models. The normal density with covariance $\Sigma$ has its level sets on the ellipsoids $\langle\mathbf x-\boldsymbol\mu,\Sigma^{-1}(\mathbf x-\boldsymbol\mu)\rangle=c$, whose axes point along the principal components.

\subsection{Positive Definite Systems and Minimisation}

The equivalence between linear systems and minimisation, which \cref{prop:quadratic-centre} established, is usually written with a factor $1/2$.

\begin{proposition}[Energy minimisation]
\label{prop:energy-minimum}
Let $A\in\mathbb R^{n\times n}$ be symmetric and positive definite, $\mathbf b\in\mathbb R^n$, and
\[
f(\mathbf x)\coloneqq\tfrac12\langle\mathbf x,A\mathbf x\rangle-\langle\mathbf b,\mathbf x\rangle .
\]
The unique minimiser of $f$ is the solution $\mathbf x^*=A^{-1}\mathbf b$ of $A\mathbf x=\mathbf b$, and for every $\mathbf x$
\[
f(\mathbf x)-f(\mathbf x^*)=\tfrac12\|\mathbf x-\mathbf x^*\|_A^2,\qquad\tfrac12\lambda_1\|\mathbf x-\mathbf x^*\|^2\le f(\mathbf x)-f(\mathbf x^*)\le\tfrac12\lambda_n\|\mathbf x-\mathbf x^*\|^2,
\]
where $\|\mathbf v\|_A\coloneqq\langle\mathbf v,A\mathbf v\rangle^{1/2}$ is the norm of the inner product $\langle\cdot,\cdot\rangle_A$.
\end{proposition}

\begin{proof}
Expanding with $A\mathbf x^*=\mathbf b$,
\[
\tfrac12\langle\mathbf x-\mathbf x^*,A(\mathbf x-\mathbf x^*)\rangle=\tfrac12\langle\mathbf x,A\mathbf x\rangle-\langle\mathbf b,\mathbf x\rangle+\tfrac12\langle\mathbf x^*,A\mathbf x^*\rangle=f(\mathbf x)-f(\mathbf x^*),
\]
since $f(\mathbf x^*)=-\langle\mathbf x^*,A\mathbf x^*\rangle/2$. The right side is positive for $\mathbf x\ne\mathbf x^*$, and the bounds are those of \cref{thm:principal-axes}.
\end{proof}

The norm $\|\cdot\|_A$ is called the \term{energy norm}, because for a stiffness matrix $A$ the quantity $\|\mathbf v\|_A^2/2$ is the potential energy of the displacement $\mathbf v$. The proposition turns the solution of a positive definite system into the minimisation of a convex function, and the minimisation of such a function into the solution of a linear system. Both directions are used: iterative methods such as the conjugate gradient method of Volume~X solve large positive definite systems by minimising $f$, and Newton's method of Volume~XIII minimises a general function by solving, at each step, the positive definite system of its local quadratic model. The level sets of $f$ are ellipsoids centred at $\mathbf x^*$, whose axes have the length ratio $\sqrt{\kappa_2(A)}$ by \cref{prop:ellipsoid}; long, thin ellipsoids, which belong to large condition numbers, are what makes such minimisation problems hard for simple descent methods, as those volumes show.

\subsection{The Second-Derivative Test Completed}

The criteria of this chapter complete the classification of critical points begun in \vcref{1}{sec:extrema-n}. Let $f\in C^2(U)$ have a critical point $\mathbf c$, and let $H\coloneqq\nabla^2f(\mathbf c)$. By \vcref{1}{thm:second-derivative-test-rn} and \cref{thm:principal-axes,thm:pd-criteria}, $f$ has a strict local minimum at $\mathbf c$ if all eigenvalues of $H$ are positive, equivalently if all leading principal minors are positive, equivalently if the Cholesky factorisation of $H$ exists; a strict local maximum if this holds for $-H$; and a saddle point if $H$ has a positive and a negative eigenvalue, which by \cref{prop:inertia-ldl} can be seen from the pivots when elimination without exchanges works. If $H$ is singular, the test is inconclusive, and by \cref{lem:psd-kernel} this is the only case in which it is: a semidefinite Hessian with $\det H\ne0$ is definite. The constant of \vcref{1}{lem:definite-bound} is $\lambda_{\min}(H)$, as \cref{sec:rayleigh} showed, and the proof of \vcref{1}{thm:second-derivative-test-rn} gives $f(\mathbf c+\mathbf h)-f(\mathbf c)\ge\lambda_{\min}(H)\|\mathbf h\|^2/4$ near $\mathbf c$. For the function $x^2+y^2+z^2+xy+yz$ of \vcref{1}{sec:extrema-n}, the Hessian has the rows $(2,1,0)$, $(1,2,1)$ and $(0,1,2)$, which is $2I+N$ for the matrix $N$ with ones next to the diagonal, a tridiagonal Toeplitz matrix; its eigenvalues are $2-\sqrt2$, $2$ and $2+\sqrt2$, those of $4I-T_3$, all positive.

\section{Symmetric Matrices in Computation: A First Look}
\label{sec:symmetric-computation}

The routines of the stardust library for symmetric and Hermitian matrices are the subject of Volume~IX, where the algorithms behind them, the reduction to tridiagonal form, the symmetric QR algorithm and divide and conquer among them, are developed. Here we test a few statements of this chapter with them. The machine and the version of the library are those described in \cref{sec:systems-practice}.

\subsection{Cholesky Against LU}

By \cref{prop:cholesky-cost}, the Cholesky factorisation needs about $n^3/3$ operations, half of the $2n^3/3$ of $PA=LU$. \Cref{tab:cholesky-timing} compares the routines \code{linalg.cholesky}, which returns the factor, and \code{linalg.lu}, which returns the factors of $PA=LU$, and the solution of a system by \code{linalg.cholesky} followed by \code{linalg.cholesky\_solve} with the general routine \code{linalg.solve}. The test matrices were $A=BB^{\mathsf T}+nI$, where $B$ has entries drawn uniformly from $[0,1)$ by the random number generator of the library with the seed $20261126$, so that $A$ is symmetric positive definite, with right-hand sides of the same kind; each time is the best of nine repetitions for $n\le2000$ and of five for larger $n$.

\begin{table}
\sdfloat{%
\begin{sdtable}{colspec={rllllll}, row{1-2}={font=\small\bfseries}}
\toprule
 & \SetCell[c=3]{l} factorisation & & & \SetCell[c=3]{l} solution of one system & & \\
\cmidrule[lr]{2-4}\cmidrule[lr]{5-7}
$n$ & \code{cholesky} & \code{lu} & ratio & Cholesky & \code{solve} & ratio \\
\midrule
$250$ & $5.9\times10^{-5}$ & $1.7\times10^{-4}$ & $2.83$ & $7.2\times10^{-5}$ & $1.6\times10^{-4}$ & $2.20$ \\
$500$ & $2.6\times10^{-4}$ & $8.2\times10^{-4}$ & $3.15$ & $3.3\times10^{-4}$ & $7.8\times10^{-4}$ & $2.40$ \\
$1000$ & $1.2\times10^{-3}$ & $2.8\times10^{-3}$ & $2.41$ & $1.3\times10^{-3}$ & $2.6\times10^{-3}$ & $1.93$ \\
$2000$ & $5.1\times10^{-3}$ & $1.5\times10^{-2}$ & $2.84$ & $5.7\times10^{-3}$ & $1.2\times10^{-2}$ & $2.17$ \\
$4000$ & $3.4\times10^{-2}$ & $8.9\times10^{-2}$ & $2.63$ & $3.6\times10^{-2}$ & $8.4\times10^{-2}$ & $2.32$ \\
$8000$ & $2.7\times10^{-1}$ & $5.0\times10^{-1}$ & $1.89$ & $2.8\times10^{-1}$ & $4.9\times10^{-1}$ & $1.77$ \\
\bottomrule
\end{sdtable}}{Times in seconds for symmetric positive definite matrices of order $n$: the Cholesky factorisation against $PA=LU$, and the solution of one system by the Cholesky factorisation and \code{linalg.cholesky\_solve} against \code{linalg.solve}, with the ratios of the times. The operation counts predict the ratio $2$ for large $n$.\label{tab:cholesky-timing}}
\end{table}

The measured ratios lie between $1.9$ and $3.2$ for the factorisations and between $1.8$ and $2.4$ for the solutions, around the predicted $2$. They do not settle on one value: the ratio for the factorisations is $2.83$ at $n=250$ and $1.89$ at $n=8000$, without a monotone trend in between, and a repetition of the whole measurement changed individual ratios by up to $0.2$. As in \cref{sec:systems-practice}, the operation count predicts the work, and how fast the work is done depends on the routine, the size and the machine; what the table shows is that the Cholesky route was faster than the general one in every case, by about the factor that the counts predict.

\subsection{When Cholesky Fails}

The Hilbert matrix $H_n$ of \cref{sec:pivoting}, with the entries $1/(i+j-1)$, is positive definite for every $n$, as \cref{sec:positive-definite} showed. Its smallest eigenvalue decreases rapidly with $n$, and in binary64 arithmetic the stored matrix, whose entries are the rounded numbers $1/(i+j-1)$, need not be positive definite. \Cref{tab:hilbert-cholesky} shows the smallest eigenvalue of $H_n$ and of the stored matrix, both computed exactly with the arbitrary-precision ball arithmetic of python-flint from the rational entries, the Frobenius norm of the difference $E$ between the two matrices, the smallest eigenvalue computed by \code{linalg.eigvalsh} in binary64, and the outcome of \code{linalg.cholesky} for the stored matrix.

\begin{table}
\sdfloat{%
\begin{sdtable}{colspec={rlllll}, row{1}={font=\small\bfseries}}
\toprule
$n$ & $\lambda_1(H_n)$ & $\lambda_1$, stored & $\|E\|_{\mathrm F}$ & \code{eigvalsh} & \code{cholesky} \\
\midrule
$10$ & $1.09\times10^{-13}$ & $1.09\times10^{-13}$ & $6.11\times10^{-17}$ & $1.09\times10^{-13}$ & succeeds \\
$11$ & $3.39\times10^{-15}$ & $3.40\times10^{-15}$ & $6.26\times10^{-17}$ & $3.40\times10^{-15}$ & succeeds \\
$12$ & $1.05\times10^{-16}$ & $1.07\times10^{-16}$ & $6.40\times10^{-17}$ & $1.09\times10^{-16}$ & succeeds \\
$13$ & $3.22\times10^{-18}$ & $8.35\times10^{-19}$ & $6.52\times10^{-17}$ & $-4.53\times10^{-18}$ & succeeds \\
$14$ & $9.88\times10^{-20}$ & $-6.32\times10^{-18}$ & $6.62\times10^{-17}$ & $-7.31\times10^{-18}$ & fails \\
$15$ & $3.02\times10^{-21}$ & $-7.13\times10^{-18}$ & $6.71\times10^{-17}$ & $-6.44\times10^{-18}$ & fails \\
$16$ & $9.20\times10^{-23}$ & $-7.16\times10^{-18}$ & $6.79\times10^{-17}$ & $-7.17\times10^{-18}$ & fails \\
\bottomrule
\end{sdtable}}{The Hilbert matrices in binary64. Second and third columns: the exact smallest eigenvalue of $H_n$ and of the stored matrix with rounded entries; fourth: the Frobenius norm of the rounding error $E$ of the entries; fifth: the smallest eigenvalue computed by \code{linalg.eigvalsh}; last: whether \code{linalg.cholesky} returned a factor or reported that the stored matrix is not positive definite.\label{tab:hilbert-cholesky}}
\end{table}

The table shows three things. First, the routine was right in every row: for $n\le13$ the stored matrix is positive definite and the factorisation succeeded, and for $n\ge14$ it is not and the routine reported this. Exact elimination in rational arithmetic on the stored matrices confirms the boundary: the first $13$ pivots are positive, and the $14$th, which is the same for every $n\ge14$ because the stored $H_{14}$ is the leading part of every larger stored matrix, is $-1.0\times10^{-13}$; by \cref{prop:cholesky-test} no Cholesky factorisation of these matrices exists. Second, the exact smallest eigenvalues of $H_n$ and of the stored matrix differ by at most $7.2\times10^{-18}$, less than $\|E\|_{\mathrm F}$, as Weyl's inequality (\cref{thm:weyl}) requires, since $\|E\|_2\le\|E\|_{\mathrm F}$; once $\lambda_1(H_n)$ is far below $\|E\|_{\mathrm F}\approx6.5\times10^{-17}$, the inequality no longer determines the sign of the smallest eigenvalue of the stored matrix, and from $n=14$ on this eigenvalue is negative. Third, the computed smallest eigenvalue for $n=13$ is negative although the stored matrix is positive definite. It differs from the exact value by $5.4\times10^{-18}$, which is small compared with $u\|H_{13}\|_2\approx2.0\times10^{-16}$; by Weyl's inequality a change of the matrix of norm $u\|H_{13}\|_2$ may move its eigenvalues by that much, so an eigenvalue computation whose accuracy is described in this way cannot decide the sign of an eigenvalue of modulus $8.4\times10^{-19}$.

\subsection{Perturbing the Eigenvalues}

\Cref{thm:weyl} says that a Hermitian perturbation $E$ moves no eigenvalue by more than $\|E\|_2$. \Cref{tab:weyl} tests this on $T_{100}$, whose eigenvalues are known exactly. For each $\delta$ we formed $20$ perturbations $E$ of each of three kinds, scaled to $\|E\|_2=\delta$: random symmetric matrices $(C+C^{\mathsf T})/2$ with entries of $C$ drawn uniformly from $[-1,1)$; random positive semidefinite matrices $\mathbf v\mathbf v^{\mathsf T}$ of rank one, with entries of $\mathbf v$ drawn in the same way; and $\mathbf q_j\mathbf q_j^{\mathsf T}$ for the eigenvectors $\mathbf q_j$, $j=1,6,11,\ldots,96$, of \cref{thm:second-difference-eigen}. We computed the eigenvalues of $T_{100}+E$ with \code{linalg.eigvalsh} and compared them, in increasing order, with the exact eigenvalues $2-2\cos(j\pi/101)$ evaluated in binary64. The table gives the largest value of $\max_k|\lambda_k(T_{100}+E)-\lambda_k(T_{100})|/\|E\|_2$ over the $20$ perturbations.

\begin{table}
\sdfloat{%
\begin{sdtable}{colspec={rlll}, row{1}={font=\small\bfseries}}
\toprule
$\delta=\|E\|_2$ & random symmetric & random rank one & $\delta\,\mathbf q_j\mathbf q_j^{\mathsf T}$ \\
\midrule
$10^{-2}$ & $0.2817$ & $0.1278$ & $1.0000$ \\
$10^{-4}$ & $0.2589$ & $0.1290$ & $1.0000$ \\
$10^{-6}$ & $0.2388$ & $0.1121$ & $1.0000$ \\
$10^{-8}$ & $0.2602$ & $0.1214$ & $1.0000$ \\
\bottomrule
\end{sdtable}}{The largest relative shift $\max_k|\lambda_k(T_{100}+E)-\lambda_k(T_{100})|/\|E\|_2$ of the eigenvalues of $T_{100}$ under symmetric perturbations of three kinds, each the largest over $20$ perturbations, computed with \code{linalg.eigvalsh}. Weyl's inequality bounds every entry by $1$.\label{tab:weyl}}
\end{table}

No entry exceeds $1$, and the perturbations along an eigenvector attain the bound, since $T_{100}+\delta\mathbf q_j\mathbf q_j^{\mathsf T}$ has the eigenvalue $\lambda_j+\delta$ and all others unchanged. The random perturbations move the eigenvalues by about a quarter of the bound, and the rank-one perturbations by about an eighth. We also checked the two-sided form of the theorem: every computed shift $\lambda_k(T_{100}+E)-\lambda_k(T_{100})$ lay between $\lambda_1(E)$ and $\lambda_n(E)$, both computed by \code{linalg.eigvalsh}, up to $10^{-14}$, and the shifts under the positive semidefinite rank-one perturbations were nonnegative up to the same tolerance, as \cref{prop:rank-one-interlacing} requires. For contrast, \cref{tab:jordan-sensitivity} showed a nonsymmetric matrix whose eigenvalue moved by $0.1$ under a perturbation of size $10^{-16}$.

\bigskip
A real quadratic form is the form of a unique symmetric matrix, and over $\mathbb C$ the forms with real values are those of Hermitian matrices; a change of variables replaces the matrix by a congruent one. The spectral theorem says that every Hermitian matrix is $Q\Lambda Q^{\mathsf H}$ with a unitary $Q$ and a real diagonal $\Lambda$: it has an orthonormal basis of eigenvectors, its spectral projections are orthogonal, the normal matrices are exactly those that unitary matrices diagonalise, and commuting Hermitian matrices are diagonalised together. In the eigenvector coordinates a form is a weighted sum of squares, which classifies conics and quadrics, gives the axes and the volume of ellipsoids, and decides definiteness by the signs of the eigenvalues. Positive definiteness is equivalent to positive eigenvalues, to a Gram representation $B^{\mathsf H}B$, to positive leading principal minors, to positive pivots, and to the Cholesky factorisation $A=R^{\mathsf H}R$, which costs $n^3/3$ operations and is the practical test; semidefiniteness needs all principal minors. Positive definiteness is preserved by congruence, inversion and Schur complements, characterises inner products on $\mathbb F^n$, and gives every positive semidefinite matrix a unique positive semidefinite square root. The Rayleigh quotient has the extreme eigenvalues as its extreme values and the eigenvectors as its critical points, its error is quadratic in the error of the vector, and $\|A\|_2=\rho(A)$ for Hermitian $A$. The Courant--Fischer theorem characterises every eigenvalue by a min--max, and it yields Weyl's inequalities, by which Hermitian eigenvalues move by at most $\|E\|_2$, and Cauchy's interlacing theorem. The law of inertia says that congruence preserves the signs of the eigenvalues, so that pivots count eigenvalues below a shift, and two forms, one of them positive definite, can be diagonalised together. Finally, the Gaussian integral of a positive definite form is $\pi^{n/2}/\sqrt{\det A}$ in every dimension, covariance matrices are positive semidefinite, and solving a positive definite system is minimising a convex quadratic function. \Cref{ch:svd} applies the spectral theorem to $A^{\mathsf H}A$ and obtains the singular value decomposition of every matrix.

\chapternotes

Hawkins~\cite{hawkins1975cauchy} traces the origins of the spectral theory of symmetric matrices in the principal axes of quadratic forms and in the problems of mechanics and astronomy of the eighteenth century, and in Cauchy's memoir of 1829, which proved that the eigenvalues of a real symmetric matrix are real and in which the interlacing of the roots of the characteristic equations of a symmetric matrix and of its principal submatrices appears; the interlacing theorem bears Cauchy's name~\cite[Sect.~4.3]{horn2013}. Sylvester stated and proved the law of inertia in 1852~\cite{sylvester1852}. The variational characterisation of the eigenvalues in between the extreme ones is due to Fischer~\cite{fischer1905}, who stated it for quadratic forms in 1905, and to Courant~\cite{courant1920}, who used it in 1920 for the eigenvalues of differential operators; Weyl's inequalities appeared in his work of 1912 on the asymptotic distribution of eigenvalues~\cite{weyl1912}. The quotient that bears Rayleigh's name is the principle by which he estimated the frequencies of vibrating systems in \emph{The Theory of Sound}~\cite{rayleigh1877}. Cholesky devised his method for the normal equations of geodesy; it was published after his death by Beno\^{\i}t in 1924~\cite{benoit1924}, and Brezinski~\cite{brezinski2006} describes his life and the manuscript of 1910 in which the method is written down. Loewner introduced the order of positive semidefinite matrices in his study of monotone matrix functions~\cite{loewner1934}, Schur proved the product theorem in 1911~\cite{schur1911}, Givens, in a report of 1954, combined the reduction of a symmetric matrix to tridiagonal form by plane rotations with the counting of eigenvalues by Sturm sequences~\cite{givens1954}, and Bunch and Kaufman gave the stable factorisation of symmetric indefinite matrices with block pivots~\cite{bunchkaufman1977}.

For the theory, Horn and Johnson~\cite{horn2013} treat Hermitian matrices, the variational principles, Weyl's inequalities, interlacing and congruence in Chapter~4 and positive definite matrices, including the Schur product theorem and the Loewner order, in Chapter~7; Axler~\cite[Chap.~7]{axler2024} proves the spectral theorem for operators without determinants, Strang~\cite{strang2016} and Meyer~\cite{meyer2000} present symmetric and positive definite matrices with many applications, and Hoffman and Kunze~\cite{hoffman1971} treat bilinear and Hermitian forms. Parlett~\cite{parlett1998} is the classical account of the symmetric eigenvalue problem, from the Rayleigh quotient and the min--max principles to the algorithms. On the computational side, Trefethen and Bau~\cite[Lecture~23]{trefethen1997} give a short introduction to the Cholesky factorisation, Higham~\cite[Chap.~10]{higham2002accuracy} analyses its rounding errors, and Golub and Van Loan~\cite[Sect.~4.2 and Chap.~8]{golub2013} treat positive definite systems and the symmetric eigenvalue problem, including bisection and the generalised problem.

\chapterreferences
